% Electronic Journal of Combinatorics format.
\documentclass[12pt]{article}
\usepackage[amsmath]{e-jc}

\usepackage{mathtools}
\usepackage{booktabs}
\usepackage{enumitem}

\MSC{11B75, 11H06, 11H31, 11Y16}
\Copyright{The author. Released under the CC BY license (International 4.0).}

% Submission copy: publication dates, volume and DOI are assigned by the journal.
% Keep the journal style file unchanged and identify this copy as a manuscript.
\makeatletter
\def\the@dateline{Manuscript, 8 September 2026\\[0.5ex]\copy@right}
\renewcommand{\ps@plain}{%
  \renewcommand{\@oddfoot}{\footnotesize Manuscript prepared for submission\hfil\thepage}}
\makeatother
\pagestyle{plain}

\newcommand{\Z}{\mathbb{Z}}
\newcommand{\Q}{\mathbb{Q}}
\newcommand{\R}{\mathbb{R}}
\newcommand{\C}{\mathbb{C}}
\newcommand{\F}{\mathbb{F}}
\newcommand{\N}{\mathbb{N}}
\newcommand{\covol}{\operatorname{covol}}
\newcommand{\rank}{\operatorname{rank}}
\newcommand{\cont}{\operatorname{cont}}
\newcommand{\adj}{\operatorname{adj}}
\newcommand{\Res}{\operatorname{Res}}
\newcommand{\rad}{\operatorname{rad}}
\newcommand{\sat}{\operatorname{sat}}
\newcommand{\Nm}{\mathrm{N}}
\newcommand{\one}{\mathbf{1}}
\newcommand{\dcube}[1]{(-#1,#1)}
\newcommand{\Lat}{\Lambda}
\newcommand{\gpoly}{g}

\title{Explicit dissociated sets from tensor lattices:\\ primitivity certificates and an effective decay rate}

\author{Zhiyu Liu}
\authortext{}{Tianjin Medical University General Hospital, Tianjin Medical University, Tianjin, China (\email{07eb1f@gmail.com}).}

\begin{document}

\maketitle

\begin{abstract}
A finite set of positive integers is \emph{dissociated} if its subset sums are pairwise distinct. We give explicit
constructions from the tensor lattices underlying the recent disproof of Erd\H{o}s's distinct subset sums conjecture
by GPT-6 Astra, as expounded by Bloom. A prescribed perturbation of a lattice in dimension $d=b^s$ has a normal vector
factoring into $s$ vectors of length $b$. We prove an exact formula for its saturation index in terms of the contents
(the gcds of the coefficients) of $s$ circulant adjugates, and characterise primitivity by polynomial gcds over prime
fields. Thus large constructions
can be certified through small matrices.

Exact integer computations give a dissociated set of size $n=43\,923$ with $\max A<0.158523\cdot2^n$, and one of size
$n=3\,841\,966$ with $\max A<0.103037\cdot2^n$, improving Bohman's explicit constant $0.22002$. Each bound extends to
every larger size by doubling and adjoining $1$.

For every odd $b\ge5$ with $3\nmid b$, we prove that the relevant first-level resultant is nonzero. The proof reduces
a hypothetical common root to a nonzero algebraic integer all of whose conjugates have absolute value less than $1$.
This nonvanishing controls every tensor depth: the scalings $Q=2^{s+1}t$ are primitive whenever $t$ is a sufficiently
large multiple of a modulus depending only on $b$. Explicit bounds on this modulus give the unconditional rate
$f(n)\le n^{-1/(50\log\log n)}$ for all sufficiently large $n$, where $f(n)$ is the infimum of $\max A/2^{n-1}$ over
dissociated sets of size $n$.
\end{abstract}

\section{Introduction}

\subsection{Background}
Let $A\subset\N$ be finite. We call $A$ \emph{dissociated} (or say that $A$ has \emph{distinct subset sums}) if the map
$S\mapsto\sum_{a\in S}a$ is injective on the subsets $S\subseteq A$. Erd\H{o}s, who called this ``perhaps my first serious
problem'', conjectured that every dissociated $A\subseteq\{1,\dots,N\}$ with $|A|=n$ satisfies $N\gg2^n$, and offered
\$500 for a proof or disproof; it is Problem~1 of the collection \cite{EP1}. Writing
\[
 f(n):=\inf\Bigl\{\frac{\max A}{2^{n-1}} : A\subset\N \text{ dissociated},\ |A|=n\Bigr\},
\]
the powers of two give $f(n)\le1$. Until 2026, every improvement of this trivial bound came from an explicit construction, and
each such bound propagates from one size to all larger sizes because $f$ is nonincreasing (Lemma~\ref{lem:mono} below). The
Conway--Guy sequence \cite{CoGu68,Gu82}, shown to consist of dissociated sets by Bohman \cite{Bo96}, gives
$N<0.23513\cdot2^n$ for $n\ge40$; Lunnon's computer search \cite{Lu88} gives $N<0.22096\cdot2^n$ for $n\ge67$; and Bohman
\cite{Bo98} obtained $f(n)\le0.44004$ for large $n$, i.e.\ $N\le0.22002\cdot2^n$, the best explicit construction prior to
the present paper.
In the other direction Erd\H{o}s and Moser \cite{Er56} proved $N\ge(\frac14-o(1))2^n/\sqrt n$; the constant was improved
several times (Elkies \cite{El86}, Aliev \cite{Al08}; see the table in \cite{St23} for the full history), the current
record $\sqrt{2/\pi}$ being due to Elkies and Gleason (unpublished) and to Dubroff, Fox and Xu \cite{DFX21}, who prove
$N\ge\binom{n}{\lfloor n/2\rfloor}$.

In September 2026 the conjecture was \emph{disproved}: for every $\varepsilon>0$ there are arbitrarily large $n$ with
$f(n)<\varepsilon$. The proof was found by GPT-6 Astra and verified in Lean \cite{Ad26}. We adopt the lattice notation
of Bloom's exposition \cite{EP1}. Its first step constructs lattices
$\Lat_s\subset V_d=\{x\in\R^d:\sum_ix_i=0\}$, $d=b^s$, which avoid the open unit cube and have small covolume.
Its second step approximates $Q\Lat_s$ by a primitive lattice $L=u^\perp\cap\Z^d$ and transfers this to dissociated
integers. Bloom invokes primitive-lattice equidistribution for that step. Horesh and Karasik \cite{HoKa23} give
quantitative error terms, but the exposition does not extract the dimension-dependent constants needed for a rate.

The original formal proof already contains a different algebraic passage to primitive lattices. The module cited
in \cite{Ad26} uses integral Smith-type diagonalisation, a saturated bidiagonal perturbation and explicit normal
weights; its perturbation buffer is obtained for each fixed matrix. It does not state a quantitative bound on
$n$ in terms of $\varepsilon$. Our contribution is a perturbation that retains the tensor coordinates: its normal
vector and saturation index reduce to $s$ calculations in dimension $b$, with an explicit perturbation constant and
local tests that yield the numerical constructions below. A new nonvanishing theorem for the first-level resultant
then provides arithmetic progressions of primitive scalings at every depth, with a modulus whose size we control
explicitly. This yields an unconditional effective decay rate for the tensor construction. The classical
monotonicity lemma below already combines with the original disproof to give qualitative convergence
$f(n)\to0$ along all $n$.

\subsection{Results}
We develop and certify the tensor-preserving perturbation just described. Throughout, $b\ge5$ is odd with $3\nmid b$, $s\ge1$, $d=b^s$,
$R_s=\Z[x_1,\dots,x_s]/(x_1^b-1,\dots,x_s^b-1)\cong\Z^d$, $\gpoly(x)=(2+x)(1-x)=2-x-x^2$, and
\begin{equation}\label{eq:Delta}
\Delta_{b,s}:=\frac{(1+2^{-b})^{(b^s-1)/(b-1)}}{(3/2)^s}
\end{equation}
is the normalised covolume $\covol(\Lat_s)/\sqrt d$ of the construction (Proposition~\ref{prop:closed}).

\paragraph{Closed form of the lattices.} We show (Proposition~\ref{prop:closed}) that
\[
\Lat_s=\bigoplus_{k=1}^{s} 2^{-k}(2+x_1)\cdots(2+x_k)(1-x_k)\,\Z[x_k,\dots,x_s],\qquad
v_s=2^{-s}\prod_{k=1}^s(2+x_k),
\]
where $\Z[x_k,\dots,x_s]\subset R_s$ is the span of the monomials not involving $x_1,\dots,x_{k-1}$.

\paragraph{An explicit primitivisation.} For $Q\in2^s\Z$ we add to each basis vector
$Q\,2^{-k}(2+x_1)\cdots(2+x_k)(1-x_k)x_k^im$ of $Q\Lat_s$ ($0\le i\le b-2$, $m$ a monomial in $x_{k+1},\dots,x_s$) the
monomial $x_1^{b-1}\cdots x_{k-1}^{b-1}x_k^im$. This defines a lattice $L_s(Q)\subset\Z^d$ of rank $d-1$.
Define \emph{level polynomials} $\varphi_k=\theta_k\gpoly+\rho_k\in\Z[x]/(x^b-1)$ recursively by $\theta_1=Q/2$, $\rho_1=1$,
$\psi_k=\adj(\overline{\varphi_k})$ (the adjugate element, $\psi_k\overline{\varphi_k}=\Nm(\varphi_k)$),
$\gamma_k=\cont(\psi_k)$, $\tilde u_k=\pm x^{b-1}\psi_k/\gamma_k$, $a_k=\langle\tilde u_k,x^{b-1}\rangle$,
$c_k=\langle\tilde u_k,2+x\rangle$, $\theta_{k+1}=2^{-(k+1)}Q\,c_1\cdots c_k$, $\rho_{k+1}=a_1\cdots a_k$
(Definition~\ref{def:levels}). The constant $K_s$ below is explicit (Proposition~\ref{prop:K}).

\begin{theorem}[normal vector and index]\label{thm:main-structure}
Let $Q\in2^s\Z$ with $Q\ge8K_s$. Then $L_s(Q)$ has rank $d-1$, its orthogonal complement is spanned by the primitive integer vector
$u=\tilde u_1\otimes\dots\otimes\tilde u_s$, i.e.\ $u(x)=\prod_k\tilde u_k(x_k)$, all coordinates of $u$ are positive, and
\[
 [\,\sat L_s(Q):L_s(Q)\,]=\prod_{k=1}^{s}\gamma_k^{\,b^{s-k}} .
\]
In particular $L_s(Q)$ is primitive if and only if $\gamma_k=\cont\adj(\overline{\varphi_k})=1$ for every $k$.
\end{theorem}

\begin{theorem}[level criterion and exceptional primes]\label{thm:levels-intro}
For $\varphi=\theta\gpoly+\rho\in\Z[x]/(x^b-1)$ with $\Nm(\varphi)\ne0$, one has $\cont\adj(\bar\varphi)=1$ if and only if
$\deg\gcd_{\F_p}(\varphi,x^b-1)\le1$ for every prime $p$. Let $P_b$ be the set of odd primes dividing $b(2^b+1)n_b$,
where $n_b=\prod_{0<i<j<b}\Res_z\bigl(\tfrac{z^b-1}{z-1},\,1+z^i+z^j\bigr)\neq0$. If $p$ is an odd prime with $p\notin P_b$
and $p\nmid\gcd(\theta,\rho)$, then $\deg\gcd_{\F_p}(\varphi,x^b-1)\le1$.
\end{theorem}

\begin{theorem}[first level, unconditional]\label{thm:A-intro}
Let $q\ge4K_1$ be an integer with $\rad(P_b)\mid q$. Then $L_1(2q)$ is primitive; the coordinates of $\tilde u_1$ are
positive integers forming a $(2q-K_1)$-dissociated set, with $\max\tilde u_1\le(2q)^{b-1}\Delta_{b,1}(1+5K_1/q)$.
Consequently, for every $l$ with $2^l\le2q-K_1$ the set $\{2^j\tilde u_{1,i}:0\le j<l,\ 0\le i<b\}$ is dissociated of size
$lb$, and $f(lb)\le\max\tilde u_1/2^{l(b-1)}$.
\end{theorem}

\paragraph{Explicit dissociated sets.} Theorem~\ref{thm:main-structure}, a perturbation lemma (Lemma~\ref{lem:cube})
and a size estimate (Lemma~\ref{lem:shape}) together give the following.

\begin{theorem}\label{thm:explicit-intro}
Let $2^s\mid Q$, $Q\ge8K_s$, and suppose $\gamma_k=1$ for $k=1,\dots,s$ (a finite computation).
Then the $d=b^s$ coordinates of $u$ are positive integers forming a $(Q-K_s)$-dissociated set, and
\[
\max u=\prod_{k=1}^s\max\tilde u_k\ \le\ Q^{d-1}\Delta_{b,s}\,(1+10K_s/Q).
\]
For every $l$ with $2^l\le Q-K_s$ the set $A=\{2^ju_i\}$ is dissociated, $|A|=ld$, and $f(ld)\le\max u/2^{l(d-1)}$.
\end{theorem}

The certificate $\gamma_1=\dots=\gamma_s=1$ requires $s$ adjugate computations of $b\times b$ integer matrices; the
integers grow with $d$, but a full $d\times d$ computation is avoided. Table~\ref{tab:main} lists the outcomes; for example
\begin{align*}
 (b,s)=(11,3):\quad & n=43\,923,\qquad\ \ \max A<0.158523\cdot2^{n},\\
 (b,s)=(13,4):\quad & n=1\,199\,562,\quad\ \max A<0.132535\cdot2^{n},\\
 (b,s)=(17,4):\quad & n=3\,841\,966,\quad\ \max A<0.103037\cdot2^{n}.
\end{align*}
These explicit dissociated sets improve Bohman's constant $0.22002$. By
Lemma~\ref{lem:mono} each bound persists for all larger sizes: there are explicit dissociated sets with
$\max A<0.158523\cdot2^{n}$ for every $n\ge43\,923$ and with $\max A<0.103037\cdot2^{n}$ for every $n\ge3\,841\,966$
(Corollary~\ref{cor:alln}).

\paragraph{Good scalings at every tensor depth.} Write $Q=2^{s+1}t$. Before contents are normalised, the level data
$\hat a_k,\hat c_k$ are polynomials in $t$ with integer coefficients (Section~\ref{sec:AP}). Let $H(b,s)$ be the
statement that $\hat a_k$ and $\hat c_k$ are coprime in $\Q[t]$ for $1\le k\le s-1$.

\begin{theorem}\label{thm:AP-intro}
For every allowed $b$ there is an integer $M_b\ge1$, independent of $s$, such that $L_s(2^{s+1}t)$ is primitive for every $t\equiv0\pmod {M_b}$ with
$2^{s+1}t\ge8K_s$. Consequently $\limsup_{l\to\infty}f(l\,b^s)\le\Delta_{b,s}$.
\end{theorem}

The key step is Theorem~\ref{thm:base-coprime}: the first-level resultant $r_b$ of \eqref{eq:base-resultant} is nonzero
for every allowed $b$. We then prove $H(b,s)$ at every depth by a homogeneous recursion
(Proposition~\ref{prop:H-reduction}), and may take $M_b=\rad(P_b)|r_b|$. Thus, for example,
$\limsup_lf(1331\,l)\le\Delta_{11,3}<0.31618$ and $\limsup_lf(2197\,l)\le\Delta_{13,3}<0.30299$ hold unconditionally.
The explicit estimate $\log M_b\le4b^3$ gives
$f(n)\le n^{-1/(20\log\log n)}$ for infinitely many $n$ (Theorem~\ref{thm:rate}), and
$f(n)\le n^{-1/(50\log\log n)}$ for all sufficiently large $n$ (Corollary~\ref{cor:rate-all}). Both statements are
proved in full in Section~\ref{sec:AP}.

\paragraph{Sizes and heights.} The classical doubling step (Lemma~\ref{lem:mono}) fills gaps between the sizes $n=ld$
without loss in the normalised maximum. Section~\ref{sec:remarks} records the exact costs of other padding and lifting
operations. We also relate decay to the height of admissible pairs: a height $h\ge b^\alpha$ and sufficiently small
excess covolume give a tensor covolume $O(d^{-\alpha})$. Transferring this to $f(n)\le n^{-\alpha+o(1)}$ requires
quantitative control of primitivisation and of the available sizes, which we state explicitly there. Vaaler's theorem
bounds the height by $O(\sqrt b)$ when the full lattice covolume is bounded (Proposition~\ref{prop:height}).

\paragraph{Organisation.} Section~\ref{sec:lattices} recasts the lattices of \cite{EP1} in the polynomial model $R_s$,
in which a cyclic shift of coordinates is multiplication by a variable, and proves the closed form.
Section~\ref{sec:perturb} contains the two perturbation lemmas and the constant $K_s$. Section~\ref{sec:tensor} defines
$L_s(Q)$ and proves Theorems~\ref{thm:main-structure}, \ref{thm:levels-intro} and \ref{thm:A-intro}.
Section~\ref{sec:main} deduces Theorems~\ref{thm:explicit-intro} and \ref{thm:AP-intro}, the monotonicity lemma with its
corollaries, and the effective decay rate. Section~\ref{sec:comp} describes the computations behind Table~\ref{tab:main}, and
Section~\ref{sec:remarks} discusses the sizes reached by the construction, the height problem and Siegel's lemma.

\subsection{Notation and conventions}\label{sec:notation}
For a lattice $L$ of rank $r$ in $\R^d$, $\covol(L)$ is the $r$-dimensional volume of a fundamental domain.
A sublattice $L\subseteq\Z^d$ is \emph{primitive} (or saturated) if $L=(L\otimes\R)\cap\Z^d$; we write
$\sat L$ for $(L\otimes\R)\cap\Z^d$. If $L\subset\Z^d$ has rank $d-1$ with basis matrix $B\in\Z^{d\times(d-1)}$, the
vector $w$ of signed maximal minors of $B$ spans $L^\perp$ and $\gcd(w)=[\sat L:L]$; a primitive lattice of rank $d-1$
is $R(u):=u^\perp\cap\Z^d$ for a primitive $u\in\Z^d$, and $\covol(R(u))=|u|_2$.
For $u\in\Z^d$ and real $Q>0$, the coordinates of $u$ are \emph{$Q$-dissociated} if the only $c\in\Z^d$ with $|c_i|<Q$
for all $i$ and $\sum_ic_iu_i=0$ is $c=0$; equivalently $R(u)\cap\dcube Q^d=\{0\}$.
The \emph{lift} of a $2^l$-dissociated vector $u$ with positive coordinates is
$A=\{2^ju_i:0\le j<l,\ 1\le i\le d\}$; two subset sums of $A$ are $\sum_ic_iu_i$ and $\sum_ic'_iu_i$ with
$0\le c_i,c'_i<2^l$, so $A$ is dissociated, $|A|=ld$, $\max A=2^{l-1}\max u$, and
\begin{equation}\label{eq:lift}
 f(ld)\le\frac{2^{l-1}\max u}{2^{ld-1}}=\frac{\max u}{2^{l(d-1)}} .
\end{equation}
$R_b=\Z[x]/(x^b-1)$ is identified with $\Z^b$ via $x^i\leftrightarrow e_i$; the standard inner product becomes
$\langle f,g\rangle=[x^0](f\bar g)$, where $\bar g(x)=g(x^{-1})$ and $[x^i]h$ denotes the coefficient of $x^i$ in $h$;
in particular $\langle f,x^i\rangle=[x^i]f$. For $a\in R_b$ let $m_a$ be the $b\times b$ matrix of multiplication by $a$
(a circulant), $\Nm(a)=\det m_a=\prod_{\zeta^b=1}a(\zeta)$, and $\adj(a)\in R_b$ the element with $m_{\adj(a)}=\adj(m_a)$,
so that $a\adj(a)=\Nm(a)$; since $\adj(m_a)$ commutes with every $m_c$, it is indeed a multiplication matrix.
Note $m_a^{T}=m_{\bar a}$. The content $\cont(\psi)$ of $\psi\in R_b$ is the gcd of its coefficients.
Finally $\one_d=(1,\dots,1)\in\R^d$, $V_d=\one_d^\perp$, and $\mu_b$ is the group of complex $b$-th roots of unity.

\section{The lattices}\label{sec:lattices}

\subsection{Admissible pairs}
Following \cite{EP1}, a pair $(\Lat,v)$ with $\Lat\subset V_d$ a lattice of rank $d-1$ and $v\in\R^d\setminus V_d$ is
\emph{admissible} if
\begin{enumerate}[label=(A\arabic*)]
\item $\Lat\cap\dcube1^d=\{0\}$;
\item $tv\in\Lat+\dcube1^d$ with $t\in\R$ implies $|t|<1$.
\end{enumerate}
Both conditions are invariant under $v\mapsto-v$; replacing $v$ by $-v$ if necessary, we always orient an admissible pair so
that its \emph{height} $h=\langle v,\one_d\rangle$ is positive. Then $\Gamma=\Lat\oplus\Z v$ satisfies
$\covol(\Gamma)=\covol(\Lat)\,h/\sqrt d$. We write $\Delta=\covol(\Lat)/\sqrt d=\covol(\Gamma)/h$.
Conditions (A1)--(A2) imply $\Gamma\cap\dcube1^d=\{0\}$, hence $\covol(\Gamma)\ge1$ by Minkowski's theorem, so
$\Delta\ge1/h$.

\medskip\noindent\emph{The base pair.} Let $b\ge3$ be odd. Following \cite{EP1}, put
\[
 \Lat_1=\tfrac12\gpoly R_b
       =\tfrac12(2+x)\bigl((1-x)R_b\bigr),
 \qquad v_1=\tfrac12(2+x).
\]
Identify $R_b$ with $\Z^b$, index coordinates by $0,\dots,b-1$, and let
$M=I+\frac12P$, where $(Py)_i=y_{i-1}$ with subscripts taken modulo $b$.
Then $(1-x)R_b=\Z^b\cap V_b$, so
$\Lat_1=M(\Z^b\cap V_b)$ and $v_1=Me_0$.

\begin{lemma}\label{lem:base-pair}
The pair $(\Lat_1,v_1)$ is admissible. Its invariants are
\[
 h_1=\frac32,\qquad
 \covol(\Gamma_1)=1+2^{-b},\qquad
 \Delta_{b,1}=\frac23(1+2^{-b}).
\]
\end{lemma}
\begin{proof}
For every $y\in\R^b$,
\[
 \lVert My\rVert_\infty
 \ge \lVert y\rVert_\infty-\tfrac12\lVert Py\rVert_\infty
 =\tfrac12\lVert y\rVert_\infty.
\]
Thus $My\in(-1,1)^b$ implies $|y_i|<2$ for every $i$.

For (A1), suppose $z\in\Z^b\cap V_b$ and $Mz\in(-1,1)^b$.
Each $z_i$ belongs to $\{-1,0,1\}$. If $z_i=1$, the inequality
$|z_i+z_{i-1}/2|<1$ forces $z_{i-1}=-1$; if $z_i=-1$, it forces
$z_{i-1}=1$. Hence a nonzero coordinate would force alternating signs
around the entire cycle, which is impossible because $b$ is odd.
Therefore $z=0$.

For (A2), suppose $t\in\R$, $z\in\Z^b\cap V_b$, and
$tv_1-Mz\in(-1,1)^b$. Set $y=te_0-z$. Then $My\in(-1,1)^b$,
the coordinates $y_1,\dots,y_{b-1}$ belong to $\{-1,0,1\}$, and
$\sum_i y_i=t$.
We first show $|y_0|<1$. If $y_0\ge1$, the coordinate-$0$ inequality
forces $y_{b-1}=-1$. Applying the same sign rule successively backwards
forces $y_1=1$, since $b-1$ is even. But then
$y_1+y_0/2\ge3/2$, contradicting the coordinate-$1$ inequality.
Replacing $y$ by $-y$ rules out $y_0\le-1$.

If one of $y_1,\dots,y_{b-1}$ is zero, every following coordinate in
this list is zero, since an integer $y_{i+1}$ satisfying
$|y_{i+1}+0/2|<1$ must vanish. The nonzero coordinates in the list
therefore form an initial alternating sequence. If this sequence is
nonempty, its first sign is opposite to that of $y_0$, by
$|y_1+y_0/2|<1$. Its sum is either $0$ or $y_1$. In the former case
$|t|=|y_0|<1$; in the latter case $t=y_0+y_1$ also lies in $(-1,1)$.
This includes the empty sequence and proves (A2) for every real $t$.

Finally, $h_1=\langle Me_0,\one_b\rangle=3/2$, and
$(\Z^b\cap V_b)\oplus\Z e_0=\Z^b$ gives $\Gamma_1=M\Z^b$.
Since $b$ is odd,
\[
 \det M=1-(-1/2)^b=1+2^{-b}.
\]
The formula for $\Delta_{b,1}$ follows by dividing by $h_1$.
\end{proof}

\subsection{Tensor products of admissible pairs}
\begin{proposition}\label{prop:tensor}
Let $(\Lat,v)$ be admissible in $\R^{b}$ and $(\Lat',v')$ admissible in $\R^{d'}$. In $\R^b\otimes\R^{d'}=\R^{bd'}$ put
\[
\Lat''=\Lat\otimes\Z^{d'}\ \oplus\ v\otimes\Lat',\qquad v''=v\otimes v'.
\]
Then $(\Lat'',v'')$ is admissible, of height $hh'$, and $\covol(\Gamma'')=\covol(\Gamma)^{d'}\covol(\Gamma')$.
\end{proposition}
\begin{proof}
View elements of $\R^b\otimes\R^{d'}$ as $b\times d'$ matrices. An element of $\Lat''$ has the form
$X=\sum_i\beta_iz_i^{T}+v\lambda^{T}$ with $\beta_i\in\Lat$, $z_i\in\Z^{d'}$, $\lambda\in\Lat'$, so its $m$-th column is
$\sum_i\beta_iz_{i,m}+\lambda_mv\in\Lat+\lambda_mv$. Coordinate sums vanish, and $\Lat\otimes\R^{d'}$, $v\otimes\R^{d'}$
intersect trivially, so $\Lat''$ is a lattice of rank $(b-1)d'+(d'-1)=bd'-1$ in $V_{bd'}$.
If $X\in\dcube1^{bd'}$, each column lies in $(\Lat+\lambda_mv)\cap\dcube1^b$, so $|\lambda_m|<1$ by (A2) for $(\Lat,v)$;
thus $\lambda\in\Lat'\cap\dcube1^{d'}=\{0\}$, and then every column lies in $\Lat\cap\dcube1^b=\{0\}$.
If $X\in\Lat''+tv''$ has all entries in $\dcube11$, the $m$-th column lies in $\Lat+(\lambda_m+tv'_m)v$, so
$|\lambda_m+tv'_m|<1$ for all $m$, i.e.\ $\lambda+tv'\in(\Lat'+tv')\cap\dcube1^{d'}$, and (A2) for $(\Lat',v')$ gives
$|t|<1$. The height is $\langle v,\one_b\rangle\langle v',\one_{d'}\rangle=hh'$. Finally
$\Gamma''=\Lat\otimes\Z^{d'}\oplus v\otimes\Gamma'$; the first summand spans $V_b\otimes\R^{d'}$ and has covolume
$\covol(\Lat)^{d'}$, while the projection of $v\otimes\Gamma'$ onto $(V_b\otimes\R^{d'})^\perp=\one_b\otimes\R^{d'}$ is
$\frac{h}{\sqrt b}\,\hat n\otimes\Gamma'$ with $\hat n=\one_b/\sqrt b$, of covolume $(h/\sqrt b)^{d'}\covol(\Gamma')$.
Hence $\covol(\Gamma'')=(\covol(\Lat)h/\sqrt b)^{d'}\covol(\Gamma')=\covol(\Gamma)^{d'}\covol(\Gamma')$.
\end{proof}

The iteration of \cite{EP1} is the same operation with the roles of the two factors exchanged; both produce admissible
pairs with identical invariants. We use the order of Proposition~\ref{prop:tensor} because it gives a clean closed form.

\begin{definition}
$(\Lat_1,v_1)$ is the base pair in the variable $x_1$; for $s\ge2$,
$(\Lat_s,v_s):=(\Lat_1,v_1)\otimes(\Lat_{s-1},v_{s-1})$, where $\Lat_{s-1}$ lives in the variables $x_2,\dots,x_s$.
\end{definition}

\begin{proposition}[closed form]\label{prop:closed}
Let $\Z[x_k,\dots,x_s]\subset R_s$ be the span of the monomials in $x_k,\dots,x_s$. Then
\[
\Lat_s=\bigoplus_{k=1}^{s}\Lat_s^{(k)},\qquad
\Lat_s^{(k)}=2^{-k}(2+x_1)\cdots(2+x_k)(1-x_k)\,\Z[x_k,\dots,x_s],
\]
$v_s=2^{-s}\prod_{k=1}^{s}(2+x_k)$, $h_s=(3/2)^s$, $\covol(\Gamma_s)=(1+2^{-b})^{(b^s-1)/(b-1)}$ and $\covol(\Lat_s)/\sqrt{b^s}=\Delta_{b,s}$ as in
\eqref{eq:Delta}. A basis of $\Lat_s^{(k)}$ consists of the vectors
$2^{-k}(2+x_1)\cdots(2+x_k)(1-x_k)\,x_k^i\,m$, $0\le i\le b-2$, $m$ a monomial in $x_{k+1},\dots,x_s$; the
denominators of $\Lat_s$ divide $2^s$.
\end{proposition}
\begin{proof}
Induction on $s$; $s=1$ is the definition of $\Lat_1$, since $(1-x)R_b$ has basis $(1-x)x^i$, $i\le b-2$.
For $s\ge2$, $\Lat_1\otimes\Z[x_2,\dots,x_s]=\frac12(2+x_1)(1-x_1)R_s=\Lat_s^{(1)}$ and
$v_1\otimes\Lat_{s-1}=\frac12(2+x_1)\bigoplus_{k\ge2}2^{-(k-1)}(2+x_2)\cdots(2+x_k)(1-x_k)\Z[x_k,\dots,x_s]
=\bigoplus_{k\ge2}\Lat_s^{(k)}$ by the induction hypothesis. The invariants follow from
Proposition~\ref{prop:tensor}: $h_s=\frac32h_{s-1}$ and
$\covol(\Gamma_s)=\covol(\Gamma_1)^{b^{s-1}}\covol(\Gamma_{s-1})$, which gives the exponent $1+b+\dots+b^{s-1}$.
\end{proof}

\section{Perturbation lemmas}\label{sec:perturb}

Let $(\Lat,v)$ be admissible in $\R^d$, let $b_1,\dots,b_{d-1}$ be a fixed basis of $\Lat$, and let
$b_1^*,\dots,b_{d-1}^*\in V_d$ be the dual basis, so that $\langle b_i,b_j^*\rangle=\delta_{ij}$. Put
\begin{equation}\label{eq:K}
 K:=\sum_{j=1}^{d-1}|b_j^*|_1 .
\end{equation}
For $\lambda=\sum_jm_jb_j\in\Lat$ we have $m_j=\langle\lambda,b_j^*\rangle$, hence
\begin{equation}\label{eq:coeff}
 \sum_j|m_j|\le K\,|\lambda|_\infty .
\end{equation}
Let $Q>0$ be real, let $r_1,\dots,r_{d-1}\in\{1,\dots,d\}$ be \emph{distinct} rows, and
\[
 L:=\operatorname{span}_\Z\{\ell_j:=Qb_j+e_{r_j}\ :\ 1\le j\le d-1\}\subset\R^d .
\]

\begin{lemma}[cube avoidance]\label{lem:cube}
If $Q>K$ then the $\ell_j$ are linearly independent and $L\cap\dcube{Q-K}^d=\{0\}$. In particular, if $L\subseteq\Z^d$
is primitive with normal vector $u$, the coordinates of $u$ are $(Q-K)$-dissociated.
\end{lemma}
\begin{proof}
Let $m\in\R^{d-1}$, $x=\sum_jm_j\ell_j=Q\lambda+\sum_jm_je_{r_j}$ with $\lambda=\sum_jm_jb_j$. Since $m_j=\langle\lambda,b_j^*\rangle$,
the bound \eqref{eq:coeff} holds for real $m$ as well, so if $m\ne0$ then $\lambda\ne0$ (the $b_j$ are a basis) and
$|x|_\infty\ge Q|\lambda|_\infty-\sum_j|m_j|\ge(Q-K)|\lambda|_\infty>0$; hence the $\ell_j$ are linearly independent over $\R$.
If moreover $m\in\Z^{d-1}\setminus\{0\}$, then $\lambda\in\Lat\setminus\{0\}$, so $|\lambda|_\infty\ge1$ by (A1) and
$|x|_\infty\ge Q-K$.
\end{proof}

\begin{lemma}[shape of the normal vector]\label{lem:shape}
\textup{(i)} Let $Q>K$ and let $u\in\R^d\setminus\{0\}$ be orthogonal to $L$. Then $\langle u,\one_d\rangle\ne0$ and,
after replacing $u$ by $-u$ if necessary, $u=\alpha(\one_d+\eta)$ with $\alpha>0$, $\eta\in V_d$ and
$|\eta|_\infty\le K/(Q-K)$. In particular, if $Q>2K$ all coordinates of $u$ have the same sign.

\textup{(ii)} If moreover $Q>2K$, $Qb_j\in\Z^d$ for all $j$, and $u\in\Z^d$ is primitive, then
\begin{gather*}
\alpha\le Q^{d-1}\Delta\,e^{K/Q}\Bigl(1+\tfrac{K^2}{(Q-K)^2}\Bigr)^{1/2},\\
\max u\le Q^{d-1}\Delta\,e^{K/Q}\Bigl(1+\tfrac{K^2}{(Q-K)^2}\Bigr)^{1/2}\Bigl(1+\tfrac{K}{Q-K}\Bigr)
\le Q^{d-1}\Delta\Bigl(1+\frac{10K}{Q}\Bigr).
\end{gather*}
\end{lemma}
\begin{proof}
(i) Write $u=\alpha\one_d+\eta'$ with $\eta'\in V_d$. From $\langle u,\ell_j\rangle=0$ and $\langle\one_d,b_j\rangle=0$
we get $Q\langle\eta',b_j\rangle+\alpha+\eta'_{r_j}=0$ for all $j$. If $\alpha=0$ then
$\eta'=\sum_j\langle\eta',b_j\rangle b_j^*=-Q^{-1}\sum_j\eta'_{r_j}b_j^*$, so $|\eta'|_\infty\le(K/Q)|\eta'|_\infty$ and
$u=0$, a contradiction. So $\alpha\ne0$; take $\alpha>0$ and $\eta=\eta'/\alpha$, so that
$\langle\eta,b_j\rangle=-(1+\eta_{r_j})/Q$ and $\eta=-Q^{-1}\sum_j(1+\eta_{r_j})b_j^*$, whence
$|\eta|_\infty\le Q^{-1}(1+|\eta|_\infty)K$, i.e.\ $|\eta|_\infty\le K/(Q-K)$, which is $<1$ when $Q>2K$.

(ii) Since $L\subseteq R(u)$, $\covol(L)\ge\covol(R(u))=|u|_2=\alpha|\one_d+\eta|_2\ge\alpha\sqrt d$.
The orthogonal projection $\pi\colon u^\perp\to V_d$ is bijective, as $\langle u,\one_d\rangle\ne0$, and it multiplies
covolumes by $\cos\vartheta=|\langle u,\one_d\rangle|/(|u|_2\sqrt d)=\sqrt d/|\one_d+\eta|_2\ge(1+|\eta|_\infty^2)^{-1/2}$.
Now $\pi(\ell_j)=Qb_j+\pi(e_{r_j})$ and $\pi(e_{r_j})=\sum_i\langle e_{r_j},b_i^*\rangle b_i=\sum_i(b_i^*)_{r_j}b_i$, so
$\pi(L)$ has basis matrix $B(QI+A)$ with $A_{ij}=(b_i^*)_{r_j}$, where $B$ is the basis matrix of $\Lat$. Hence
\begin{align*}
\covol(\pi(L))&=Q^{d-1}\covol(\Lat)\,|\det(I+A/Q)|\\
&\le Q^{d-1}\sqrt d\,\Delta\prod_i\Bigl(1+\frac{|b_i^*|_1}{Q}\Bigr)\le Q^{d-1}\sqrt d\,\Delta\,e^{K/Q} .
\end{align*}
Here the first inequality is Hadamard's, using $\sum_j|A_{ij}|\le|b_i^*|_1$ (the rows $r_j$ are distinct), and the
second follows from \eqref{eq:K}. Therefore $\alpha\sqrt d\le\covol(L)=\covol(\pi(L))/\cos\vartheta\le Q^{d-1}\sqrt d\,\Delta\,e^{K/Q}(1+|\eta|_\infty^2)^{1/2}$.
The last inequality of the lemma is elementary for $x=K/Q\le\frac12$: $e^x\le1+1.3x$,
$(1+x^2/(1-x)^2)^{1/2}\le1+x$, $1+x/(1-x)\le1+2x$, and $(1+1.3x)(1+x)(1+2x)\le1+10x$.
\end{proof}

\begin{proposition}[the constant $K_s$]\label{prop:K}
Let $K_1=\sum_{i\le b-2}|b_i^*|_1$ for the base basis $b_i=\frac12\gpoly x^i$ of $\Lat_1$, and let
$c_i=\langle b_i^*,v_1\rangle/h_1$. For the basis of $\Lat_s$ in Proposition~\ref{prop:closed}, which is
$\{b_i\otimes m\}\cup\{v_1\otimes b'_j\}$ with $b'_j$ the basis of $\Lat_{s-1}$ and $d'=b^{s-1}$, a dual basis is
\[
(b_i\otimes m)^*=b_i^*\otimes m-c_i\,\one_b\otimes\bigl(m-\tfrac1{d'}\one_{d'}\bigr),\qquad
(v_1\otimes b'_j)^*=\tfrac{1}{h_1}\one_b\otimes(b'_j)^*.
\]
Consequently the constant \eqref{eq:K} for this basis of $\Lat_s$ equals
\[
K_s=d'\sum_{i}\Bigl(\bigl|b_i^*-c_i\bigl(1-\tfrac1{d'}\bigr)\one_b\bigr|_1+|c_i|\,b\,\tfrac{d'-1}{d'}\Bigr)+\tfrac{b}{h_1}K_{s-1}
\ \le\ d'\bigl(1+\tfrac{4b}{3}\bigr)K_1+\tfrac{2b}{3}K_{s-1}.
\]
\end{proposition}
\begin{proof}
Direct verification of $\langle\cdot,\cdot\rangle=\delta$, using $\langle\one_b,b_i\rangle=0$,
$\langle\one_b,v_1\rangle=h_1$, $\langle m-\frac1{d'}\one_{d'},b'_j\rangle=(b'_j)_m$, and that all vectors lie in
$V_{bd'}$. The $\ell^1$ norms are read off blockwise; the crude bound uses
$|c_i|\le|b_i^*|_1|v_1|_\infty/h_1=\frac23|b_i^*|_1$.
\end{proof}
Numerically $K_1(5)=424/55\approx7.71$, $K_1(7)\approx14.19$, $K_1(11)\approx32.20$, $K_1(13)\approx44.03$, and,
for instance, $K_2(5)\approx98.25$, $K_3(11)\approx1.78\cdot10^4$, $K_4(13)\approx5.52\cdot10^5$. We checked the
recursion against a direct computation of the dual basis for $(b,s)\in\{(5,2),(7,2),(5,3)\}$.


\section{The tensor perturbation}\label{sec:tensor}

Fix $b\ge5$ odd with $3\nmid b$, $s\ge1$, $d=b^s$, and $R_s$, $x_1,\dots,x_s$ as in Proposition~\ref{prop:closed}.

\begin{definition}[the lattice $L_s(Q,\rho)$]\label{def:L}
Let $Q\in2^s\Z_{>0}$ and $\rho\in\Z_{>0}$. For $1\le k\le s$ put
\[
\Phi_k=\Phi_k^{(Q,\rho)}:=\frac{Q}{2^k}(1-x_k)\prod_{j\le k}(2+x_j)\ +\ \rho\prod_{j<k}x_j^{\,b-1}\ \in R_s,
\]
and let $L_s(Q,\rho)$ be the $\Z$-span of the $d-1$ vectors
$\ell_{k,i,m}:=x_k^i\,m\,\Phi_k$, where $1\le k\le s$, $0\le i\le b-2$, and $m$ runs over the monomials in
$x_{k+1},\dots,x_s$. We write $L_s(Q)=L_s(Q,1)$.
\end{definition}
By Proposition~\ref{prop:closed}, $\ell_{k,i,m}=Q\,\beta_{k,i,m}+\rho\,\mu_{k,i,m}$, where the $\beta_{k,i,m}$ form the
basis of $\Lat_s$ and $\mu_{k,i,m}=x_1^{b-1}\cdots x_{k-1}^{b-1}x_k^im$ are \emph{distinct} monomials (the variable
$x_k$ is the first one with exponent $\le b-2$), none of them equal to $\prod_jx_j^{b-1}$. Thus $L_s(Q)$ is exactly the
lattice of Section~\ref{sec:perturb} for $\Lat_s$ with the tensor basis and distinct rows, and
$L_s(Q,\rho)=\rho\cdot\operatorname{span}\{(Q/\rho)\beta_{k,i,m}+\mu_{k,i,m}\}$ is the same construction with the real
scale $Q/\rho$.

\begin{definition}[level data]\label{def:levels}
Set $\theta_1=Q/2$ and $\rho_1=\rho$. For $1\le k\le s$, provided $\theta_k,\rho_k>0$, put
\begin{gather*}
\varphi_k=\theta_k\gpoly+\rho_k\in R_b,\qquad \Nm_k=\Nm(\varphi_k),\qquad
\psi_k=\adj(\overline{\varphi_k})\quad(\text{so }\psi_k\overline{\varphi_k}=\Nm_k),\\
\gamma_k=\cont(\psi_k),\qquad \tilde u_k=\varepsilon_k\,x^{b-1}\psi_k/\gamma_k\quad
(\varepsilon_k\in\{\pm1\}\text{ with }\langle\tilde u_k,\one_b\rangle>0),\\
a_k=\langle\tilde u_k,x^{b-1}\rangle=\varepsilon_k\psi_{k,0}/\gamma_k,\qquad
c_k=\langle\tilde u_k,2+x\rangle=2\tilde u_{k,0}+\tilde u_{k,1},\\
\theta_{k+1}=\frac{Q}{2^{k+1}}\,c_1\cdots c_k,\qquad \rho_{k+1}=\rho\,a_1\cdots a_k\qquad(k<s),
\end{gather*}
where $\tilde u_{k,i}$ and $\psi_{k,i}$ denote the coefficients of $x^i$; thus $\tilde u_{k,i}=\varepsilon_k\psi_{k,i+1}/\gamma_k$
with indices mod $b$. (Since $2^s\mid Q$, the $\theta_k$ are integers for $k\le s$; the recursion is not continued beyond $k=s$.)
\end{definition}

\begin{lemma}[non-degeneracy]\label{lem:nondeg}
If $Q\ge8K_s\rho$, then for every $k\le s$: $\theta_k,\rho_k>0$, $\Nm_k\ne0$, and $\tilde u_k$ has all coordinates
positive (so $a_k,c_k>0$). Moreover $2\theta_k/\rho_k\ge Q/\rho$.
\end{lemma}
\begin{proof}
For $\theta,\rho'>0$ the element $\theta\gpoly+\rho'$ is not a zero divisor of $R_b\otimes\Q$: a root $\zeta\in\mu_b$ would
satisfy $\gpoly(\zeta)=-\rho'/\theta\in\Q_{<0}$, but $\gpoly(\zeta)=2-\zeta-\zeta^2$ is rational only for
$\zeta\in\{1,-1,\omega,\bar\omega\}$ ($\omega$ a primitive cube root of unity), and
$\mu_b\cap\{1,-1,\omega,\bar\omega\}=\{1\}$ because $b$ is odd and $3\nmid b$, while $\gpoly(1)=0$.
Now induct on $k$. Given $\theta_k,\rho_k>0$ we have $\Nm_k\ne0$ and $\psi_k\ne0$. The lattice
$L^{(k)}:=\operatorname{span}\{x^i\varphi_k:i\le b-2\}=\rho_k\operatorname{span}\{(2\theta_k/\rho_k)b_i+e_i\}$ is of the type
of Lemma~\ref{lem:shape} with $K=K_1$ and real scale $Q_k':=2\theta_k/\rho_k$, and its orthogonal complement is spanned
by $x^{b-1}\psi_k$ (Proposition~\ref{prop:level} below). If $Q_k'\ge8K_1$, Lemma~\ref{lem:shape}(i) shows that
$\tilde u_k$ has coordinates of one sign, all in $[\alpha(1-\frac17),\alpha(1+\frac17)]$, whence
$c_k/a_k\ge3\cdot\frac{6/7}{8/7}>2$. Since $K_s\ge(2b/3)^{s-1}K_1\ge K_1$ by Proposition~\ref{prop:K}, we have
$Q'_1=Q/\rho\ge8K_1$, and inductively $2\theta_{k+1}/\rho_{k+1}=(Q/2^{k}\rho)\prod_{j\le k}(c_j/a_j)\ge Q/\rho\ge8K_1$.
\end{proof}

\begin{proposition}[one level]\label{prop:level}
Let $\varphi\in R_b$ with $\Nm(\varphi)\ne0$ and $L=\operatorname{span}\{x^i\varphi:0\le i\le b-2\}\subset\Z^b$. Then the
vector of signed maximal minors of $L$ is $\pm x^{b-1}\adj(\bar\varphi)$; hence $L^\perp$ is spanned by
$x^{b-1}\adj(\bar\varphi)$ and $[\sat L:L]=\cont\adj(\bar\varphi)$. Moreover
\begin{align*}
\cont\adj(\bar\varphi)=1&\iff \rank_{\F_p}(m_\varphi)\ge b-1\ \text{ for all primes }p\\
&\iff \deg\gcd_{\F_p}(\varphi,x^b-1)\le1\ \text{ for all primes }p .
\end{align*}
\end{proposition}
\begin{proof}
The columns of $m_\varphi$ are $x^j\varphi$, $0\le j\le b-1$, so the minors of $L$ (delete row $i$) are the cofactors
$C_{i,b-1}$ of $m_\varphi$; i.e.\ $w=\pm(\text{row }b-1\text{ of }\adj m_\varphi)=\pm\adj(m_\varphi)^{T}e_{b-1}
=\pm\adj(m_{\bar\varphi})e_{b-1}=\pm x^{b-1}\adj(\bar\varphi)$, using $m_\varphi^T=m_{\bar\varphi}$ and the fact that
$\adj(m_{\bar\varphi})$ is multiplication by $\adj(\bar\varphi)$. The content of $w$ is the index.
For the equivalences: $\adj(M)\equiv0\pmod p$ iff all $(b-1)$-minors of $M$ vanish mod $p$ iff $\rank_{\F_p}M\le b-2$;
and $\rank_{\F_p}m_\varphi=\rank_{\F_p}m_{\bar\varphi}=\dim_{\F_p}\bigl(\varphi\cdot\F_p[x]/(x^b-1)\bigr)=b-\deg\gcd_{\F_p}(\varphi,x^b-1)$,
since the image of multiplication by $\varphi$ in the principal ideal ring $\F_p[x]/(x^b-1)$ is the ideal generated by
$\gcd(\varphi,x^b-1)$.
\end{proof}

\begin{theorem}[normal vector]\label{thm:normal}
Let $Q\ge8K_s\rho$ and let the level data be as in Definition~\ref{def:levels}. Then $L_s(Q,\rho)$ has rank $d-1$ and
$L_s(Q,\rho)\subseteq R(u)$, where
\[
u:=\tilde u_1\otimes\dots\otimes\tilde u_s,\qquad u(x_1,\dots,x_s)=\prod_{k=1}^s\tilde u_k(x_k),
\]
is a primitive vector of $\Z^{d}$ with positive coordinates.
\end{theorem}
\begin{proof}
Rank $d-1$: Lemma~\ref{lem:cube} for the real scale $Q/\rho>K_s$. Primitivity and positivity of $u$: the coordinates of
$u$ are the products $\prod_k\tilde u_{k,i_k}$, positive by Lemma~\ref{lem:nondeg}, and the content of a tensor product
of integer vectors is the product of the contents (Gauss's lemma), so $\cont u=1$.

Orthogonality. Fix $(k,i,m)$. Since
$\langle u,\ell_{k,i,m}\rangle=[x^0]\bigl(u\,\overline{\Phi_k}\,x_k^{-i}\bar m\bigr)=[x_k^im]\bigl(u\overline{\Phi_k}\bigr)$,
it suffices to show that every monomial of $u\overline{\Phi_k}$ whose $x_1,\dots,x_{k-1}$-exponents vanish has
$x_k$-exponent $b-1$. Now
\begin{align*}
u\overline{\Phi_k}={}&\frac{Q}{2^k}\prod_{j<k}\Bigl[\tilde u_j(x_j)(2+x_j^{-1})\Bigr]\cdot\tilde u_k(x_k)\bar\gpoly(x_k)\cdot\prod_{j>k}\tilde u_j(x_j)\\
&+\rho\prod_{j<k}\Bigl[\tilde u_j(x_j)x_j^{-(b-1)}\Bigr]\cdot\tilde u_k(x_k)\cdot\prod_{j>k}\tilde u_j(x_j),
\end{align*}
a sum of two products of univariate factors. Extracting, for each $j<k$, the part with $x_j$-exponent $0$ amounts to
replacing the $j$-th factor by its constant coefficient: $[x_j^0]\tilde u_j(2+x_j^{-1})=2\tilde u_{j,0}+\tilde u_{j,1}=c_j$ and
$[x_j^0]\tilde u_jx_j^{-(b-1)}=[x_j^0]\tilde u_jx_j=\tilde u_{j,b-1}=a_j$. Hence the relevant part of $u\overline{\Phi_k}$ is
\[
\bigl(\theta_k\bar\gpoly(x_k)+\rho_k\bigr)\tilde u_k(x_k)\prod_{j>k}\tilde u_j
=\overline{\varphi_k}\,\tilde u_k\prod_{j>k}\tilde u_j
=\pm\frac{\Nm_k}{\gamma_k}\,x_k^{b-1}\prod_{j>k}\tilde u_j(x_j),
\]
by $\overline{\varphi_k}\psi_k=\Nm_k$. All its monomials have $x_k$-exponent $b-1$, as required.
\end{proof}

\begin{theorem}[index formula]\label{thm:index}
Under the hypotheses of Theorem~\ref{thm:normal},
\[
[\,R(u):L_s(Q,\rho)\,]=\prod_{k=1}^{s}\gamma_k^{\,b^{s-k}} .
\]
In particular $L_s(Q,\rho)$ is primitive if and only if $\gamma_1=\dots=\gamma_s=1$.
\end{theorem}
\begin{proof}
Induction on $s$. For $s=1$, $L_1(Q,\rho)=\operatorname{span}\{x^i\varphi_1\}$ and the index is $\gamma_1$ by
Proposition~\ref{prop:level}. Let $s\ge2$, $d'=b^{s-1}$, and view $\Z^d=\Z^b\otimes\Z^{d'}$ with $x_1$ in the first factor.
The family $k=1$ spans $L^{(1)}\otimes\Z^{d'}$ with $L^{(1)}=\operatorname{span}\{x_1^i\varphi_1\}$, whose orthogonal complement
in $\R^b$ is $\R\tilde u_1$. For $k\ge2$ write $\Phi_k=(2+x_1)\otimes A_k+x_1^{b-1}\otimes B_k$ with
$A_k=\frac{Q}{2^k}(1-x_k)\prod_{2\le j\le k}(2+x_j)$ and $B_k=\rho\prod_{2\le j<k}x_j^{b-1}$ in $R_{s-1}(x_2,\dots,x_s)$.
Decompose $\R^b\otimes\R^{d'}=W\oplus W^\perp$ with $W=\tilde u_1^\perp\otimes\R^{d'}$ and
$W^\perp=\R\tilde u_1\otimes\R^{d'}$, and identify $W^\perp$ with $\R^{d'}$ via $\tilde u_1/|\tilde u_1|_2\otimes y\mapsto y$.
The projection of $\ell_{k,i,m}$ ($k\ge2$) onto $W^\perp$ is
\[
\frac{1}{|\tilde u_1|_2}\bigl(\langle\tilde u_1,2+x_1\rangle A_k+\langle\tilde u_1,x_1^{b-1}\rangle B_k\bigr)x_k^im
=\frac{1}{|\tilde u_1|_2}\,(c_1A_k+a_1B_k)\,x_k^im ,
\]
and $c_1A_k+a_1B_k=\frac{c_1Q/2}{2^{k-1}}(1-x_k)\prod_{2\le j\le k}(2+x_j)+a_1\rho\prod_{2\le j<k}x_j^{b-1}$, so these
projections form, up to the factor $1/|\tilde u_1|_2$, the basis of $L':=L_{s-1}(Q',\rho')$ with $Q'=c_1Q/2$ and
$\rho'=a_1\rho$ in the variables $x_2,\dots,x_s$. The level data of $L'$ are $(\theta_2,\rho_2),\dots,(\theta_s,\rho_s)$:
indeed $Q'/2=\theta_2$, $\rho'=\rho_2$, and the recursion is the same. By Lemma~\ref{lem:nondeg},
$Q'/\rho'=(c_1/a_1)(Q/2\rho)\ge Q/\rho\ge8K_s\ge8K_{s-1}$, so the induction hypothesis applies to $L'$: its normal vector is
$u'=\tilde u_2\otimes\dots\otimes\tilde u_s$ and $[R(u'):L']=\prod_{k\ge2}\gamma_k^{b^{s-k}}$.
Since $L'$ has rank $d'-1$, the projection onto $W^\perp$ is injective on $L_{\ge2}:=\operatorname{span}\{\ell_{k,i,m}:k\ge2\}$;
hence $L_s(Q,\rho)\cap W=L^{(1)}\otimes\Z^{d'}$, the projection of $L_s(Q,\rho)$ onto $W^\perp$ is $L'/|\tilde u_1|_2$, and
\[
\covol L_s(Q,\rho)=\covol(L^{(1)}\otimes\Z^{d'})\cdot\covol\bigl(L'/|\tilde u_1|_2\bigr)
=\covol(L^{(1)})^{d'}\cdot\frac{\covol L'}{|\tilde u_1|_2^{\,d'-1}} ,
\]
where $\covol(L^{(1)}\otimes\Z^{d'})=\covol(L^{(1)})^{d'}$ because $L^{(1)}\otimes\Z^{d'}$ is the orthogonal direct sum
of the $d'$ copies $L^{(1)}\otimes e_m$. As $u=\tilde u_1\otimes u'$ is primitive, $\covol R(u)=|u|_2=|\tilde u_1|_2|u'|_2$, so
\[
[R(u):L_s(Q,\rho)]=\frac{\covol L_s(Q,\rho)}{\covol R(u)}=\Bigl(\frac{\covol L^{(1)}}{|\tilde u_1|_2}\Bigr)^{d'}\frac{\covol L'}{|u'|_2}
=\gamma_1^{\,d'}\prod_{k\ge2}\gamma_k^{\,b^{s-k}} .\qedhere
\]
\end{proof}

\begin{proposition}[exceptional primes]\label{prop:badprimes}
Let $n_b=\prod_{0<i<j<b}\Res_z\bigl(\frac{z^b-1}{z-1},1+z^i+z^j\bigr)$ and let $P_b$ be the set of odd primes dividing
$b(2^b+1)n_b$. Then $n_b\ne0$, so $P_b$ is finite, and for $\varphi=\theta\gpoly+\rho$ with $\theta,\rho\in\Z$ and a prime $p$:
\begin{enumerate}[label=(\alph*)]
\item if $p\mid\gcd(\theta,\rho)$ then $\deg\gcd_{\F_p}(\varphi,x^b-1)=b$;
\item if $p\mid\theta$ and $p\nmid\rho$ then $\deg\gcd_{\F_p}(\varphi,x^b-1)=0$;
\item if $p=2$ and $\rho$ is odd then $\deg\gcd_{\F_2}(\varphi,x^b-1)=0$;
\item if $p$ is odd, $p\notin P_b$ and $p\nmid\gcd(\theta,\rho)$, then $\deg\gcd_{\F_p}(\varphi,x^b-1)\le1$.
\end{enumerate}
\end{proposition}
\begin{proof}
For roots of unity, $1+\zeta_1+\zeta_2=0$ forces $\{\zeta_1,\zeta_2\}=\{\omega,\bar\omega\}$, which is impossible in $\mu_b$
because $3\nmid b$; hence every factor of $n_b$ is a nonzero integer. (a) and (b) are clear: $\varphi\equiv0$, respectively a
nonzero constant, mod $p$. (c): if $\theta$ is even, $\varphi\equiv\rho\equiv1$; if $\theta$ is odd,
$\varphi\equiv\gpoly+1\equiv1+x+x^2\pmod2$, whose roots are the primitive cube roots of unity, which are not $b$-th roots of
unity because $3\nmid b$.
(d): Let $\F=\overline{\F}_p$. Since $p\nmid b$, $x^b-1$ is separable over $\F$; fix a prime $\mathfrak p\mid p$ of
$\Z[\zeta_b]$; reduction modulo $\mathfrak p$ maps $\mu_b$ bijectively onto the roots of $x^b-1$ in $\F$.
As $x^b-1$ is squarefree over $\F$, $\deg\gcd_{\F_p}(\varphi,x^b-1)$ is the number of distinct common roots, so suppose
$\zeta_1\ne\zeta_2$ are common roots. If $p\mid\theta$ then $p\nmid\rho$ and $\varphi\equiv\rho$ has no roots. So $p\nmid\theta$
and $\gpoly(\zeta_1)=\gpoly(\zeta_2)=-\rho/\theta$; since $\gpoly(a)-\gpoly(c)=-(a-c)(1+a+c)$, we get $1+\zeta_1+\zeta_2=0$ in $\F$.
If $\zeta_1,\zeta_2\ne1$, lifting to $\tilde\zeta_i\in\mu_b$ gives $\mathfrak p\mid(1+\tilde\zeta_1+\tilde\zeta_2)$; writing
$\tilde\zeta_1=\eta^{i}$, $\tilde\zeta_2=\eta^{j}$ for a primitive $\eta\in\mu_b$, the algebraic integer $1+\eta^i+\eta^j$
divides $\Res_z(\frac{z^b-1}{z-1},1+z^i+z^j)$, so $p\mid n_b$, a contradiction. If $\zeta_1=1$, then $\zeta_2=-2$, so
$(-2)^b=1$ in $\F$, i.e.\ $p\mid2^b+1$, again a contradiction.
\end{proof}

\begin{theorem}[first level]\label{thm:A}
Let $b\ge5$ be odd, $3\nmid b$, and let $q\ge4K_1$ be an integer with $\rad(P_b)\mid q$. Then $L_1(2q)$ is primitive.
Its normal vector $\tilde u_1=\pm x^{b-1}\adj(\bar\varphi_1)$, $\varphi_1=q\gpoly+1$, has positive coordinates forming a
$(2q-K_1)$-dissociated set with $\max\tilde u_1\le(2q)^{b-1}\Delta_{b,1}(1+5K_1/q)$, and for $2^l\le2q-K_1$ the lift
$\{2^j\tilde u_{1,i}\}$ is a dissociated set of size $lb$ with $f(lb)\le\max\tilde u_1/2^{l(b-1)}$.
\end{theorem}
\begin{proof}
By Proposition~\ref{prop:level} it suffices to check $\deg\gcd_{\F_p}(\varphi_1,x^b-1)\le1$ for all $p$. For odd $p\mid q$
use (b); for odd $p\nmid q$ we have $p\notin P_b$ and $\gcd(q,1)=1$, so (d) applies; $p=2$ is (c). The remaining assertions
are Lemmas~\ref{lem:cube} and \ref{lem:shape} with $K=K_1$, $Q=2q$, together with \eqref{eq:lift}.
\end{proof}

\begin{example}\label{ex:5}
Let $b=5$ and $Q=16$, so $\varphi_1=8\gpoly+1=17-8x-8x^2$ and $\overline{\varphi_1}=17-8x^3-8x^4$. One finds
\[
\psi_1=\adj(\overline{\varphi_1})=62017+48704x+49792x^2+53704x^3+52104x^4,
\]
$\Nm_1=266321=11^2\cdot31\cdot71$, with $\psi_1\overline{\varphi_1}=266321$ in $R_5$, $\gamma_1=1$, and
$\tilde u_1=x^4\psi_1=(48704,\,49792,\,53704,\,52104,\,62017)$. Here $K_1=424/55$, so Lemma~\ref{lem:cube} guarantees
that the coordinates are $8$-dissociated (they are in fact $16$-dissociated but not $17$-dissociated, by brute force), and
$\max\tilde u_1/Q^4=62017/65536\approx0.946$, compared with $\Delta_{5,1}=11/16$; the ratio tends to $\Delta_{5,1}$ as
$Q\to\infty$, e.g.\ it is $0.6906$ for $Q=8200$. The lift with $l=3$ is the dissociated $15$-element set
$\{\tilde u_{1,i},2\tilde u_{1,i},4\tilde u_{1,i}\}$. Finally $a_1=62017$ and $c_1=2\cdot48704+49792=147200$ are coprime,
so the second level can be built on top of this one.
\end{example}

\section{Explicit dissociated sets}\label{sec:main}

\begin{theorem}\label{thm:explicit}
Let $b\ge5$ be odd with $3\nmid b$, $s\ge1$, $d=b^s$, $Q\in2^s\Z$ with $Q\ge8K_s$, and suppose that the level data of
Definition~\ref{def:levels} (with $\rho=1$) satisfy $\gamma_1=\dots=\gamma_s=1$. Then $u=\tilde u_1\otimes\dots\otimes\tilde u_s$
is a vector of $d$ positive integers whose coordinates are $(Q-K_s)$-dissociated, and
\[
\max u=\prod_{k=1}^s\max_i\tilde u_{k,i}\ \le\ Q^{d-1}\Delta_{b,s}\Bigl(1+\frac{10K_s}{Q}\Bigr).
\]
For every $l\ge1$ with $2^l\le Q-K_s$, the set $A=\{2^ju_i:0\le j<l,\ 1\le i\le d\}$ is dissociated, $|A|=ld$, and
\[
f(ld)\ \le\ \frac{\max u}{2^{l(d-1)}}\ \le\ \Delta_{b,s}\Bigl(1+\frac{10K_s}{Q}\Bigr)\Bigl(\frac{Q}{2^l}\Bigr)^{d-1}.
\]
\end{theorem}
\begin{proof}
Theorems~\ref{thm:normal} and \ref{thm:index} give $L_s(Q)=R(u)$; Lemma~\ref{lem:cube} with $K=K_s$
(Proposition~\ref{prop:K}) gives $R(u)\cap\dcube{Q-K_s}^d=\{0\}$; Lemma~\ref{lem:shape}(ii) bounds $\max u$; the product
formula holds because all $\tilde u_{k,i}>0$; finally apply \eqref{eq:lift}.
\end{proof}

The sizes $ld$ in Theorem~\ref{thm:explicit} form a sparse set. The following classical observation, used by Conway and
Guy and by Bohman \cite{Bo98} to pass from one size to all larger sizes, fills the gaps at no cost.

\begin{lemma}[monotonicity]\label{lem:mono}
If $A$ is dissociated with $|A|=n\ge1$, then $2A\cup\{1\}$ is dissociated with $n+1$ elements and maximum $2\max A$.
Consequently $f(n+1)\le f(n)$ for all $n\ge1$, and for every $m\ge n$ the explicit set
$A_m=2^{m-n}A\cup\{2^j:0\le j<m-n\}$ is dissociated with $|A_m|=m$ and $\max A_m/2^{m}=\max A/2^{n}$
(the added set is empty when $m=n$).
\end{lemma}
\begin{proof}
The elements of $2A$ are even and $1$ is odd, so $2A\cup\{1\}$ has $n+1$ elements. If two subsets of $2A\cup\{1\}$ have
the same sum, the two sums have the same parity, so either both subsets contain $1$ or neither does; removing $1$ and
halving gives two subsets of $A$ with the same sum, hence the same subsets. The maximum is $2\max A$ because $\max A\ge1$,
and $2\max A/2^{n}=\max A/2^{n-1}$; taking the infimum over $A$ gives $f(n+1)\le f(n)$, and iterating the construction
$m-n$ times gives $A_m$.
\end{proof}

\begin{corollary}\label{cor:alln}
Table~\ref{tab:main} and Lemma~\ref{lem:mono} give explicit dissociated sets $A$ of every size $n$ with
\begin{align*}
\max A&<0.158523\cdot2^{n}&&\text{for every } n\ge43\,923,\\
\max A&<0.132535\cdot2^{n}&&\text{for every } n\ge1\,199\,562,\\
\max A&<0.103037\cdot2^{n}&&\text{for every } n\ge3\,841\,966.
\end{align*}
In particular $f(n)\le0.206073$ for all $n\ge3\,841\,966$.
\end{corollary}
\begin{proof}
The rows $(b,s)=(11,3)$, $(13,4)$ and $(17,4)$ of Table~\ref{tab:main} provide dissociated sets of sizes $43\,923$,
$1\,199\,562$ and $3\,841\,966$ with $\max A/2^{|A|}$ below $0.158523$, $0.132535$ and $0.103037$ respectively
(Theorem~\ref{thm:explicit} with the certificates of Section~\ref{sec:comp}); apply Lemma~\ref{lem:mono}.
\end{proof}

\subsection{An arithmetic progression of good scalings}\label{sec:AP}
Write $Q=2^{s+1}t$, where $t\in\Z_{>0}$. Run the recursion of Definition~\ref{def:levels}
formally, without dividing by contents and without sign normalisation:
\begin{gather*}
\hat\theta_1=2^st,\qquad \hat\rho_1=1,\qquad
\hat\varphi_k=\hat\theta_k\gpoly+\hat\rho_k,\qquad
\hat\psi_k=\adj(\overline{\hat\varphi_k}),\\
\hat a_k=[x^0]\hat\psi_k,\qquad
\hat c_k=2[x^1]\hat\psi_k+[x^2]\hat\psi_k,\\
\hat\theta_{k+1}=2^{s-k}t\prod_{j\le k}\hat c_j,\qquad
\hat\rho_{k+1}=\prod_{j\le k}\hat a_j\qquad(1\le k<s).
\end{gather*}
All these quantities lie in $\Z[t]$. Equivalently,
\begin{equation}\label{eq:formal-step}
\hat\theta_{k+1}=\tfrac12\hat\theta_k\hat c_k,\qquad
\hat\rho_{k+1}=\hat\rho_k\hat a_k.
\end{equation}
The product formula above shows integrality despite the factor $1/2$.
If the first $k-1$ actual levels are primitive, then
$(\theta_k,\rho_k)=\sigma_k(\hat\theta_k(t),\hat\rho_k(t))$ for a common sign
$\sigma_k\in\{1,-1\}$. Indeed the signs chosen in Definition~\ref{def:levels} enter
both products in the same way, and
$\adj(-\varphi)=(-1)^{b-1}\adj(\varphi)=\adj(\varphi)$ since $b$ is odd.
If level $k$ is primitive as well, then
$(a_k,c_k)=\varepsilon_k(\hat a_k(t),\hat c_k(t))$ for another common sign.

\begin{definition}[the coprimality condition $H(b,s)$]\label{def:H}
We say that $H(b,s)$ holds if, for every $1\le k<s$, the polynomials
$\hat a_k,\hat c_k$ are coprime in $\Q[t]$; equivalently,
$R_k:=\Res_t(\hat a_k,\hat c_k)\ne0$ for every $1\le k<s$.
For $s=1$ the condition is empty.
\end{definition}

The condition can in fact be checked at the first level alone. To see this, define
homogeneous polynomials $\mathcal A_b,\mathcal C_b\in\Z[\theta,\rho]$ by
\begin{align*}
\mathcal A_b(\theta,\rho)
&=[x^0]\adj(\theta\bar\gpoly+\rho),\\
\mathcal C_b(\theta,\rho)
&=2[x^1]\adj(\theta\bar\gpoly+\rho)+[x^2]\adj(\theta\bar\gpoly+\rho).
\end{align*}
Both have degree $b-1$. Put $\kappa_b=(2^b+1)/3\in\Z$.
Since multiplication by $\bar\gpoly$ has kernel $\R\one_b$, its adjugate matrix is
$\kappa_b\one_b\one_b^T$: the product of its nonzero eigenvalues is
$b(2^b+1)/3$. Equivalently,
$\adj(\bar\gpoly)=\kappa_b\sum_{i=0}^{b-1}x^i$ in $R_b$.
Together with the adjugate of a scalar matrix, this gives the polynomial identities
\begin{equation}\label{eq:homogeneous-boundaries}
\begin{aligned}
\mathcal A_b(0,\rho)&=\rho^{b-1},&
\mathcal C_b(0,\rho)&=0,\\
\mathcal A_b(\theta,0)&=\kappa_b\theta^{b-1},&
\mathcal C_b(\theta,0)&=(2^b+1)\theta^{b-1}.
\end{aligned}
\end{equation}
These identities hold over $\Z$, so they remain valid after reduction modulo a prime.
Let
\begin{equation}\label{eq:base-resultant}
r_b:=\Res_q\bigl(\mathcal A_b(q,1),\mathcal C_b(q,1)\bigr)\in\Z.
\end{equation}

\begin{theorem}[nonvanishing of the base resultant]\label{thm:base-coprime}
For every odd $b\ge5$ with $3\nmid b$, the polynomials
$\mathcal A_b(q,1)$ and $\mathcal C_b(q,1)$ have no common complex root.
Consequently $r_b\ne0$ for every such $b$.
\end{theorem}
\begin{proof}
Put $\varphi=1+q\gpoly$ and write
\[
 \psi=\adj(\bar\varphi)=\sum_{i=0}^{b-1}\psi_i x^i.
\]
At $q=0$ we have $\psi=1$, so $\mathcal A_b(0,1)=1$.
Suppose henceforth that $q\ne0$ and that
\begin{equation}\label{eq:common-root-data}
 \psi_0=0,\qquad 2\psi_1+\psi_2=0.
\end{equation}

First, $\psi\ne0$, including when $\Nm(\varphi)=0$.
Indeed, the values $\gpoly(\zeta)$ for $\zeta\in\mu_b$ are distinct.
If $\zeta\ne\eta$ and $\gpoly(\zeta)=\gpoly(\eta)$, then
\[
 0=-(\zeta-\eta)(1+\zeta+\eta),
\]
so $\zeta+\eta=-1$. Two unit complex numbers with this sum are the
primitive cube roots of unity, contrary to $3\nmid b$.
The eigenvalues of multiplication by $\bar\varphi$ are
$1+q\gpoly(\zeta^{-1})$, $\zeta\in\mu_b$; thus at most one is zero.
This matrix has rank at least $b-1$, and its adjugate is a nonzero
circulant. Its first column, the coefficient vector of $\psi$, is
therefore nonzero.

Set $\lambda=(1+2q)/q$, and extend the coefficients $\psi_i$
periodically with period $b$. Extracting coefficients of $x^i$ in
$\bar\varphi\psi=\Nm(\varphi)$ gives
\begin{equation}\label{eq:common-root-recurrence}
 \psi_{i+2}=\lambda\psi_i-\psi_{i+1}
 \qquad(1\le i\le b-1).
\end{equation}
If $\psi_1=0$, then \eqref{eq:common-root-data} gives $\psi_2=0$,
and the recurrence forces every coefficient to vanish, a
contradiction. We may therefore put $y_i=\psi_i/\psi_1$. Then
\[
 y_1=1,\qquad y_2=-2,\qquad y_b=0,\qquad y_{b+1}=1.
\]

Let $r,w$ be the roots of $z^2+z-\lambda=0$, so $r+w=-1$.
If $r=w=-1/2$, the unique sequence satisfying
\eqref{eq:common-root-recurrence} and the given first two values is
\[
 y_i=(3i-2)(-1/2)^{i-1}\qquad(1\le i\le b+1).
\]
In particular $y_b\ne0$, a contradiction. If $r\ne w$, the same
recurrence gives
\[
 y_i=\frac{(r-1)r^{i-1}-(w-1)w^{i-1}}{r-w}
 \qquad(1\le i\le b+1).
\]
The equalities $y_b=0$ and $y_{b+1}=1$ imply, respectively, that
$(r-1)r^{b-1}=(w-1)w^{b-1}$ and that their common value is $1$.
Hence
\begin{equation}\label{eq:common-root-units}
 r^{b-1}(r-1)=1,\qquad w^{b-1}(w-1)=1,\qquad r+w=-1.
\end{equation}

These three identities are impossible. To see this, $r$ is a
nonzero algebraic integer, being a root of the monic polynomial
$z^b-z^{b-1}-1$, and $w=-1-r$ belongs to $\Q(r)$.
For any embedding $\sigma:\Q(r)\hookrightarrow\C$, write
$z=\sigma(r)$ and $v=\sigma(w)=-1-z$. The identities
\eqref{eq:common-root-units} give
\[
 |z|^{b-1}|z-1|=|v|^{b-1}|v-1|=1.
\]
Writing $a=\operatorname{Re}z$, we also have
\[
 |v|^2-|z|^2=2a+1,\qquad
 |v-1|^2-|z-1|^2=3(2a+1).
\]
If $a>-1/2$, both factors in the product for $v$ are strictly
larger than the corresponding factors for $z$, which is
impossible; if $a<-1/2$, the reverse inequalities give the same
contradiction. Thus $a=-1/2$. In particular $|z-1|\ge3/2$, so
\[
 |\sigma(r)|^{b-1}\le\frac23<1
\]
for every embedding $\sigma$. The norm of $r$ is a nonzero integer,
but the product of the absolute values of all its conjugates is
strictly less than $1$, a final contradiction. Therefore
\eqref{eq:common-root-data} cannot hold, proving the theorem.
\end{proof}

\begin{samepage}
\begin{proposition}[coprimality at every depth]\label{prop:H-reduction}
For every odd $b\ge5$ with $3\nmid b$, the condition $H(b,s)$ holds for every $s\ge1$.
\end{proposition}
\end{samepage}
\begin{proof}
By Theorem~\ref{thm:base-coprime}, $r_b\ne0$. The only common zero of $\mathcal A_b$ and $\mathcal C_b$ in $\C^2$
is $(0,0)$. Indeed, if $\rho\ne0$, homogeneity would give a common zero
$q=\theta/\rho$ of the two polynomials in \eqref{eq:base-resultant};
if $\rho=0$ and $\theta\ne0$, then
$\mathcal A_b(\theta,0)=\kappa_b\theta^{b-1}\ne0$.

Consider the map
\begin{equation}\label{eq:homogeneous-map}
F_b(\theta,\rho)
=\bigl(\tfrac12\theta\mathcal C_b(\theta,\rho),\,
        \rho\mathcal A_b(\theta,\rho)\bigr).
\end{equation}
It sends every nonzero pair in $\C^2$ to a nonzero pair.
When $\theta=0$ and $\rho\ne0$, its second coordinate is $\rho^b\ne0$;
when $\rho=0$ and $\theta\ne0$, its first coordinate is
$\tfrac12(2^b+1)\theta^b\ne0$.
When both coordinates of the input are nonzero, a zero output would require
$\mathcal A_b=\mathcal C_b=0$, which has just been excluded.

For any $z\in\C$, the initial pair $(2^sz,1)$ is nonzero, and
\eqref{eq:formal-step} iterates $F_b$.
Thus, at every level,
\[
(\hat\theta_k(z),\hat\rho_k(z))\ne(0,0),\qquad
(\hat a_k(z),\hat c_k(z))\ne(0,0).
\]
The polynomials $\hat a_k,\hat c_k$ have no common complex root and are therefore
coprime in $\Q[t]$. This proves $H(b,s)$ for every $s$.
\end{proof}

We next make the finite set of congruence obstructions independent of $s$.
As before, $\rad(P_b)$ denotes the product of the primes in the finite set $P_b$.
By Theorem~\ref{thm:base-coprime}, the positive integer
\begin{equation}\label{eq:uniform-modulus}
M_b:=\rad(P_b)\,|r_b|
\end{equation}
is defined for every allowed $b$.

\begin{lemma}[local avoidance]\label{lem:local-avoidance}
Let $p$ be an odd prime such that $p\notin P_b$ and $p\nmid r_b$.
Then $\mathcal A_b,\mathcal C_b$ have no common nonzero zero in
$\overline{\F}_p^2$, and the map $F_b$ of \eqref{eq:homogeneous-map}, reduced
modulo $p$, sends every nonzero pair to a nonzero pair.
\end{lemma}
\begin{proof}
The Sylvester adjugate identity gives polynomials $U,V\in\Z[q]$ with
\[
U(q)\mathcal A_b(q,1)+V(q)\mathcal C_b(q,1)=r_b.
\]
After reduction modulo $p$, the right side is nonzero. Thus the two univariate
polynomials have no common root, even if their degrees decrease after reduction.
Homogeneity excludes a common zero with $\rho\ne0$.
For $\rho=0$ and $\theta\ne0$, the value
$\mathcal A_b(\theta,0)=\kappa_b\theta^{b-1}$ is nonzero: every prime divisor
of $\kappa_b$ divides $2^b+1$ and is in $P_b$.

The three cases in the proof of Proposition~\ref{prop:H-reduction} now apply
over $\overline{\F}_p$. At $\theta=0$, the second coordinate of $F_b$ is
$\rho^b$; at $\rho=0$, the first is $\tfrac12(2^b+1)\theta^b$.
Both boundary values are nonzero when the input is nonzero, because
$p\nmid2(2^b+1)$. With $\theta\rho\ne0$, a zero output would require a common
zero of $\mathcal A_b,\mathcal C_b$.
\end{proof}

\begin{samepage}
\begin{theorem}\label{thm:AP}
Let $b\ge5$ be odd with $3\nmid b$, let $s\ge1$, and let $M_b$ be as in \eqref{eq:uniform-modulus}.
If $t\equiv0\pmod {M_b}$ and $Q=2^{s+1}t\ge8K_s$, then
$\gamma_1=\dots=\gamma_s=1$; hence $L_s(Q)$ is primitive and
Theorem~\ref{thm:explicit} applies.
The same modulus $M_b$ works for every $s\ge1$.
\end{theorem}
\end{samepage}
\begin{proof}
For $s=1$, write $Q=4t=2q$. Then $q=2t\ge4K_1$ and $\rad(P_b)\mid q$,
so Theorem~\ref{thm:A} applies. Suppose henceforth that $s\ge2$.

We first record properties of the formal recursion modulo each prime.
If $p\mid2t$, then
$\hat\theta_j=2^{s+1-j}t\prod_{i<j}\hat c_i$ is divisible by $p$ for every
$1\le j\le s$; for $p=2$ this uses $s+1-j\ge1$.
Starting with $\hat\rho_1=1$, \eqref{eq:homogeneous-boundaries} gives
\[
\hat a_j\equiv\hat\rho_j^{b-1}\not\equiv0\pmod p,
\qquad
\hat\rho_{j+1}=\hat\rho_j\hat a_j\not\equiv0\pmod p.
\]
Thus $\hat\rho_j$ is nonzero modulo $p$ at every level. This argument does not
divide by $2$ in characteristic $2$.

If $p\nmid2t$, then $p$ is odd, $p\notin P_b$ and $p\nmid r_b$, because
$M_b\mid t$. The initial pair $(2^st,1)$ is nonzero modulo $p$.
Lemma~\ref{lem:local-avoidance} and \eqref{eq:formal-step} show that
$(\hat\theta_j(t),\hat\rho_j(t))$ is nonzero modulo $p$ for every $j$.

Now induct on the actual level $k$, assuming $\gamma_j=1$ for $j<k$.
The observation preceding Definition~\ref{def:H} gives
$(\theta_k,\rho_k)=\pm(\hat\theta_k(t),\hat\rho_k(t))$ with a common sign.
If $p\mid2t$, then $p\mid\theta_k$ and $p\nmid\rho_k$, so
Proposition~\ref{prop:badprimes}(b) gives degree zero for the common gcd.
For every remaining prime, $p\notin P_b$ and
$p\nmid\gcd(\theta_k,\rho_k)$, so Proposition~\ref{prop:badprimes}(d)
gives degree at most one. The level data are nondegenerate by
Lemma~\ref{lem:nondeg}, since $Q\ge8K_s$.
Proposition~\ref{prop:level} therefore gives $\gamma_k=1$, completing the induction.
\end{proof}

\begin{corollary}\label{cor:limsup}
Let $b\ge5$ be odd with $3\nmid b$, and let $s\ge1$. Then
\[
\limsup_{l\to\infty}f(l\,b^s)\le\Delta_{b,s}.
\]
\end{corollary}
\begin{proof}
Fix $b,s$ and the modulus $M_b$ of Theorem~\ref{thm:AP}. Put $W=2^{s+1}M_b$.
For each sufficiently large $l$, choose
\[
Q=W\left\lceil\frac{2^l+K_s}{W}\right\rceil,
\qquad 2^l+K_s\le Q<2^l+K_s+W.
\]
Then $Q\ge8K_s$ and Theorems~\ref{thm:AP} and \ref{thm:explicit} give
\[
f(l\,b^s)\le\Delta_{b,s}(1+10K_s/Q)(Q/2^l)^{b^s-1}
\longrightarrow\Delta_{b,s}.
\]
\end{proof}

\begin{corollary}\label{cor:numbers}
In particular,
\[
\limsup_{l\to\infty}f(1331\,l)\le\Delta_{11,3}<0.31618,\qquad
\limsup_{l\to\infty}f(2197\,l)\le\Delta_{13,3}<0.30299.
\]
Since $f$ is nonincreasing, $\lim_{n\to\infty}f(n)\le\Delta_{13,3}$; hence
$f(n)<0.303$ for all sufficiently large $n$.
\end{corollary}
\begin{proof}
Apply Corollary~\ref{cor:limsup} and Lemma~\ref{lem:mono}.
The displayed numerical bounds follow from \eqref{eq:Delta}.
\end{proof}
Corollary~\ref{cor:numbers} requires only numerical evaluation of \eqref{eq:Delta}.
The explicit sets of
Table~\ref{tab:main} give the sharper effective statement
$f(n)\le0.206073$ for all $n\ge3\,841\,966$ (Corollary~\ref{cor:alln}).

We give quantitative bounds before deriving a rate. In what follows, $\log$
denotes the natural logarithm, while $\log_2$ denotes the logarithm to base $2$.

\begin{lemma}[height and perturbation bounds]\label{lem:AP-heights}
For every allowed $b$, the modulus in \eqref{eq:uniform-modulus} satisfies
\[
\log M_b\le4b^3.
\]
For every $s\ge1$ and $d=b^s$, the constants of Proposition~\ref{prop:K} satisfy
\[
K_1\le\frac{2(b^2-1)}3,\qquad
K_s\le5K_1d\le\frac{10}{3}b^2d.
\]
\end{lemma}
\begin{proof}
For an integer polynomial $P$, let $\|P\|_{\mathrm{coeff},1}$ be the sum of the
absolute values of its coefficients. This norm is submultiplicative.
Every entry in the matrix of multiplication by $q\bar\gpoly+1$ has this norm
at most $3$. Expanding a $(b-1)\times(b-1)$ minor as a determinant gives
\[
\|\mathcal A_b(q,1)\|_{\mathrm{coeff},1},\quad
\|\mathcal C_b(q,1)\|_{\mathrm{coeff},1}
\le S_b:=3(b-1)!3^{b-1}.
\]
The extra factor $3$ allows for the linear combination defining $\mathcal C_b$.
Both polynomials have degree $b-1$, since their leading coefficients are
$\kappa_b$ and $2^b+1$, respectively, by \eqref{eq:homogeneous-boundaries}.
Every row of their Sylvester matrix has Euclidean norm at most $S_b$.
Hadamard's inequality therefore gives
\begin{equation}\label{eq:base-resultant-height}
\log|r_b|\le2(b-1)\log S_b\le2b^2\log(3b).
\end{equation}

Each factor in the definition of $n_b$ is a product, over $b-1$ complex roots
of unity, of numbers of absolute value at most $3$. Since these resultants
are nonzero integers by Proposition~\ref{prop:badprimes},
\[
\log\rad(P_b)\le\log b+\log(2^b+1)
+\binom{b-1}{2}(b-1)\log3.
\]
Combining this with \eqref{eq:base-resultant-height} gives
\[
\log M_b\le\log b+\log(2^b+1)
+\binom{b-1}{2}(b-1)\log3+2b^2\log(3b)\le4b^3
\]
for $b\ge5$. For the last inequality, one may use
$\log b\le b$, $\log(2^b+1)\le b$, $\log(3b)\le b$ and $\log3<2$.

For the bound on $K_1$, let $D_i\in V_b$ be the dual vectors to
$e_i-e_{i+1}$, $0\le i\le b-2$. Directly,
\[
\|D_i\|_1=\frac{2(i+1)(b-i-1)}b,
\qquad \sum_{i=0}^{b-2}\|D_i\|_1=\frac{b^2-1}{3}.
\]
The base basis is the image of this chain basis under $M=I+P/2$, where $P$ is
the cyclic shift. This map and its inverse preserve $V_b$, so
$b_i^*=M^{-T}D_i$. The convergent geometric series for $M^{-T}$ gives
$\|M^{-T}\|_{1\to1}\le\sum_{j\ge0}2^{-j}=2$, proving the estimate for $K_1$.
Finally, write $E_s=K_s/b^s$. The recurrence in Proposition~\ref{prop:K} gives
\[
E_s\le\left(\frac1b+\frac43\right)K_1+\frac23E_{s-1},\qquad E_1=K_1/b.
\]
Induction yields $E_s\le(4+3/b)K_1\le5K_1$, which proves the result.
\end{proof}

\begin{theorem}[effective decay rate]\label{thm:rate}
$f(n)\le n^{-c/\log\log n}$ for infinitely many $n$, with $c=1/20$.
\end{theorem}
\begin{proof}
For each sufficiently large integer $s$, let $b=b(s)$ be the least odd integer
not divisible by $3$ which is at least $2s\log_2s$.
The gaps between consecutive positive integers congruent to $1$ or $5$ modulo
$6$ are at most $4$, so
\[
2s\log_2s\le b\le2s\log_2s+4\le s^2
\]
for all sufficiently large $s$. Consequently
$2^{-b}b^{s-1}\le s^{-2s}s^{2s-2}=s^{-2}\le1/2$.
With $d=b^s$, this implies
\[
\log\bigl(\Delta_{b,s}(3/2)^s\bigr)
=\frac{b^s-1}{b-1}\log(1+2^{-b})
\le\frac{b}{b-1}\,2^{-b}b^{s-1}<1,
\]
and hence $\Delta_{b,s}\le e(2/3)^s$.

Put $W=2^{s+1}M_b$ and choose
\begin{equation}\label{eq:rate-choice}
l=\left\lceil\log_2(W+8K_s)\right\rceil
+\left\lceil\log_2(200d)\right\rceil,\qquad
Q=W\left\lceil\frac{2^l+K_s}{W}\right\rceil.
\end{equation}
Thus $2^l\ge200d(W+8K_s)$ and
$2^l+K_s\le Q<2^l+K_s+W$. In particular, $Q\ge8K_s$ and $Q/2^{s+1}$ is a
multiple of $M_b$. Theorem~\ref{thm:AP} certifies all levels as primitive.
The complete loss, including rounding to the progression, satisfies
\begin{align*}
\log\left((Q/2^l)^{d-1}(1+10K_s/Q)\right)
&\le\frac{(d-1)(K_s+W)+10K_s}{2^l}\\
&\le\frac{d(W+8K_s)}{2^l}\le\frac1{200}.
\end{align*}
Here the middle inequality holds for $d\ge2$.
Theorem~\ref{thm:explicit} therefore produces a set of size $n_s=ld$ with
\begin{equation}\label{eq:rate-sequence}
f(n_s)\le e^{1+1/200}(2/3)^s<3(2/3)^s.
\end{equation}

Lemma~\ref{lem:AP-heights} gives
$\log_2W\le s+1+4b^3/\log2$ and
$\log_2K_s\le\log_2(10/3)+2\log_2b+\log_2d$.
It follows directly from \eqref{eq:rate-choice} that
\[
l=O(b^3+s+\log d).
\]
Since $b\le s^2$ and $\log d=s\log b\le2s\log s$, this upper bound is a fixed
polynomial in $s$, whereas $d\ge5^s$.
In particular, for all sufficiently large $s$,
\begin{equation}\label{eq:rate-size}
d\le n_s=ld\le d^4.
\end{equation}
All constants here are absolute, and the displayed estimates give effective
upper bounds for the parameter choices.

As $\log b\le2\log s\le2\log\log d$, we have
$s\ge\log d/(2\log\log d)$.
Combining \eqref{eq:rate-sequence} and \eqref{eq:rate-size} yields
\[
f(n_s)\le3\exp\left(-\frac{\log(3/2)\log d}{2\log\log d}\right)
\le3\,n_s^{-\log(3/2)/(8\log\log n_s)}
\le n_s^{-1/(20\log\log n_s)}
\]
for all sufficiently large $s$, because $\log(3/2)/8>1/20$.
Finally, $n_s\ge5^s$ tends to infinity, so these inequalities supply infinitely
many distinct sizes.
\end{proof}

\begin{corollary}\label{cor:rate-all}
$f(n)\le n^{-c'/\log\log n}$ for all sufficiently large $n$, with $c'=1/50$.
\end{corollary}
\begin{proof}
For each sufficiently large $s$, use the size $n_s$ produced above.
Given large $n$, choose $s$ maximal with $n_s\le n$.
Such an $s$ exists and is finite, because $n_s\ge5^s$; no monotonicity of the
sequence $(n_s)$ is needed. By maximality,
$n<n_{s+1}\le d_{s+1}^4$. Hence, for all sufficiently large $s$,
\[
\log n<4(s+1)\log b(s+1)
\le8(s+1)\log(s+1)\le16s\log s.
\]
On the other hand, $n\ge n_s\ge5^s$ implies
$\log\log n\ge\log s+\log\log5>\log s$.
Thus $s>\log n/(16\log\log n)$.
Lemma~\ref{lem:mono} and \eqref{eq:rate-sequence} give
\[
f(n)\le f(n_s)<3(2/3)^s
\le3\exp\left(-\frac{\log(3/2)\log n}{16\log\log n}\right)
\le n^{-1/(50\log\log n)}
\]
for all sufficiently large $n$, since $\log(3/2)/16>1/50$.
\end{proof}

The progression in Theorem~\ref{thm:AP} supplies the specific gap bound $W$
used in \eqref{eq:rate-choice}. This controls the distance from a good scaling
to $2^l+K_s$, which is the quantity needed for the lifting estimate.
The constants $1/20$ and $1/50$ retain substantial slack in the size estimates.

\section{Computations and certificates}\label{sec:comp}

The algebraic certificates use exact integers and rational numbers in Python, with \texttt{sympy} for symbolic
polynomial and matrix calculations. Floating-point quantities are used only for exploratory checks and displays
explicitly marked as approximations. The table generator computes circulant adjugates by the Cayley--Hamilton
identity; small cases are compared with fraction-free Bareiss elimination and direct symbolic adjugates.

Every one of the fifteen rows has a compressed exact witness. For each level the witness stores $\theta,\rho$, the
determinant, and the first two coefficients of $\psi=\adj(\bar\varphi)$, from which every other coefficient is exactly
reconstructible. Writing $a=2\theta+\rho$, the independent verifier works backwards through
\[
a\psi_i=\theta(\psi_{i+1}+\psi_{i+2})\qquad(1\le i<b),
\]
with subscripts modulo $b$, and checks every division and the cyclic closure. It independently computes
\[
T_0=2,\quad T_1=-\theta,\quad T_j=-\theta T_{j-1}+a\theta T_{j-2},\qquad
\Nm(\varphi)=a^b+T_b-\theta^b.
\]
Indeed, if $r_1,r_2$ are the roots of $\theta z^2+\theta z-a$, then
$T_j=\theta^j(r_1^j+r_2^j)$, and taking the product over the $b$th roots of unity gives this determinant formula for
odd $b$. Finally the verifier checks $a\psi_0-\theta(\psi_1+\psi_2)=\Nm(\varphi)\ne0$; thus it fixes the scale of the
adjugate as well as its direction. It then reconstructs the tensor maximum, checks every content and the theorem's
inequalities, and verifies the decimal upper bounds. Full SHA-256 digests identify the witness files and the
reconstructed integers. The data are in \texttt{results/certificates.json} and \texttt{results/witnesses/}; run
\texttt{python tools/verify\_certificates.py} from the repository root to check all rows independently of the
production adjugate implementation. The repository is linked in the Acknowledgements.

\begin{enumerate}[leftmargin=*]
\item \emph{Correctness checks.} For $(b,s)\in\{(5,1),(5,2),(3,2),(7,2)\}$ and several $Q$ we built the basis of $L_s(Q)$
explicitly, computed its vector of maximal minors directly, and verified that it equals
$\pm\prod_k\gamma_k^{b^{s-k}}\cdot u$ with $u=\otimes_k\tilde u_k$, as predicted by Theorems~\ref{thm:normal} and
\ref{thm:index} (for instance $(b,s)=(3,2)$, $Q=12$ gives $\gamma=(19,169)$ and index $19^3\cdot169$; the case $3\mid b$
is not covered by our theorems but the algebra of Section~\ref{sec:tensor} still applies). The invariants $h_s$,
$\covol(\Gamma_s)$ and $\Delta_{b,s}$ of Proposition~\ref{prop:closed} were verified numerically. For $(b,s)=(3,2)$,
enumeration of $7^8=5\,764\,801$ coefficient vectors found no nonzero point violating (A1) in the chosen finite box,
and no violation of (A2) at the three tested real parameters $t=1,5/4,3/2$ in that box. These finite checks do not
replace the proofs of Lemma~\ref{lem:base-pair} and Proposition~\ref{prop:tensor}, which apply to every real $t$.
\item \emph{Brute-force dissociation.} The vectors of Definition~\ref{def:levels}, lifted with $b=5$, $l\in\{3,4,5\}$ (for
$Q=16,24,40$) and with $b=7$, $l=3$ (for $Q=24$), of sizes $15$, $20$, $25$ and $21$, as well as the $25$-element vectors
$u=\tilde u_1\otimes\tilde u_2\in\Z^{25}$ of Definition~\ref{def:levels} for $(b,s)=(5,2)$ and $Q\in\{108,116,124\}$
(scalings below the threshold $8K_2$ of Theorem~\ref{thm:explicit}, which therefore does not cover them), were verified to be dissociated by a
meet-in-the-middle enumeration of all signed sums. These are independent checks of small examples; the one-level
examples also fail the sufficient divisibility or size hypotheses of Theorem~\ref{thm:A}, so that theorem is not
being invoked to certify them.
\item \emph{Sharpness of Lemma~\ref{lem:cube}.} For $b=5$ and $Q\in\{16,24,40\}$ the coordinates of $\tilde u_1$ are in fact
exactly $Q$-dissociated, whereas the lemma guarantees $Q-K_1\approx Q-7.7$. We do not know whether
$L_s(Q)\cap\dcube Q^d=\{0\}$ holds in general; if it did, one could take $Q=2^l$ whenever $2^l$ is a good scaling, and the
factor $(Q/2^l)^{d-1}$ in Theorem~\ref{thm:explicit} would disappear.
\item \emph{Level criterion.} For $b\in\{5,7,9,11,13\}$ and $8\le q<40$ the criterion of Proposition~\ref{prop:level},
evaluated by factoring $\Res(\varphi_1,x^b-1)$ and computing polynomial gcds mod $p$, agreed with the directly computed index
in all $160$ cases. For $3\mid b$ the first level is never primitive: $\varphi_1(\omega)=3q+1$ for both primitive cube
roots of unity $\omega$, so every prime factor of $3q+1$ is exceptional; this is why we assume $3\nmid b$.
\item \emph{Theorem~\ref{thm:A}.} $P_5=\{3,5,11\}$, $P_7=\{3,7,43\}$, $P_{11}=\{3,11,23,683\}$; for $q=m\cdot\rad(P_b)$,
$1\le m\le40$, the first level was primitive in all $120$ cases. The level-one polynomials satisfy
$\hat a_1(q)=\Nm(q)-\frac qb\Nm'(q)$ with $\Nm(q)=\prod_{\zeta\ne1}(1+q\gpoly(\zeta))$, and
$\Res_q(\hat a_1,\hat c_1)$ equals $2\cdot11\cdot19$, $2^3\cdot7\cdot43\cdot379$,
$2\cdot3^5\cdot5^2\cdot23^3\cdot683\cdot14797$ and $3^3\cdot13\cdot53^3\cdot61\cdot2731\cdot626592619$ for $b=5,7,11,13$;
the primes $19$ (for $b=5$) and $7$ (for $b=7$) are exactly the second-level failures observed at $Q=20$ and $Q=12$.
\item \emph{Checks of first-level coprimality.} The four nonzero resultants in the preceding item give numerical
checks of Theorem~\ref{thm:base-coprime}, which proves nonvanishing for every allowed $b$ without computation.
The file \texttt{results/first\_level\_resultants.json} contains the complete integer coefficients of the first-level
polynomials and adjugates, their determinants, and the resultant values and factorisations. The accompanying
\texttt{exp09} script checks the polynomial adjugate identity and the resultants with exact arithmetic.
The separate \texttt{exp10} script checks the circulant and recurrence identities and computes polynomial gcds for
all odd $3\le b\le101$, including $3\mid b$ as a control. These calculations supplement the general proof.
\item \emph{The table.} For each $(b,s)$ in Table~\ref{tab:main} we chose
$l=\lceil\log_2(200(d-1)K_s)\rceil$ as an initial scale, and the least $t$ with
$Q=2^{s+1}t\ge2^l+K_s$ and $\gamma_1=\dots=\gamma_s=1$; at most three values of $t$ were needed in every case.
Rounding and search are accounted for using the actual value of $Q$. Put $z=(d-1)(Q-2^l)/2^l$.
The verifier checks the rational inequalities $0\le z<1$ and $(1-z)^{-1}<1.006$ in every row. The binomial theorem gives
\[
(Q/2^l)^{d-1}\le\sum_{j\ge0}z^j=(1-z)^{-1}<1.006.
\]
The exact maximum is also checked against the shape estimate in every row.
For $(17,4)$ the four adjugates involve integers of about $3\cdot10^6$ bits.
The two bound columns round the exact rationals $\max u/2^{l(d-1)}$ and $\max u/2^{l(d-1)+1}$ upwards to six decimal
places. For a positive rational $p/q$, the printed bound is
$\lceil10^6p/q\rceil/10^6$, computed entirely with integers. The $\Delta_{b,s}$ column is an approximation to five
decimal places, not an asserted upper bound.
\end{enumerate}

\begin{table}[ht]
\centering\small
\begin{tabular}{rrrrrrccc}
\toprule
$b$ & $s$ & $d=b^s$ & $l$ & $n=ld$ & $Q$ & $f(n)\le$ & $\max A/2^n\le$ & $\Delta_{b,s}$\\
\midrule
5 & 1 & 5 & 13 & 65 & 8200 & 0.690649 & 0.345325 & 0.68750\\
7 & 1 & 7 & 15 & 105 & 32784 & 0.674015 & 0.337008 & 0.67188\\
11 & 1 & 11 & 16 & 176 & 65572 & 0.670804 & 0.335402 & 0.66699\\
13 & 1 & 13 & 17 & 221 & 131120 & 0.669767 & 0.334884 & 0.66675\\
5 & 2 & 25 & 19 & 475 & 524392 & 0.537132 & 0.268566 & 0.53457\\
7 & 2 & 49 & 22 & 1078 & 4194584 & 0.474515 & 0.237258 & 0.47299\\
11 & 2 & 121 & 25 & 3025 & 33555528 & 0.448813 & 0.224407 & 0.44706\\
13 & 2 & 169 & 26 & 4394 & 67110656 & 0.447207 & 0.223604 & 0.44520\\
5 & 3 & 125 & 25 & 3125 & 33555168 & 0.771243 & 0.385622 & 0.76915\\
7 & 3 & 343 & 28 & 9604 & 268438368 & 0.463429 & 0.231715 & 0.46171\\
11 & 3 & 1331 & 33 & 43923 & 8589952400 & 0.317045 & 0.158523 & 0.31617\\
13 & 3 & 2197 & 34 & 74698 & 17179903936 & 0.304339 & 0.152170 & 0.30299\\
17 & 3 & 4913 & 37 & 181781 & 137439055184 & 0.298073 & 0.149037 & 0.29699\\
13 & 4 & 28561 & 42 & 1199562 & 4398047063424 & 0.265070 & 0.132535 & 0.26412\\
17 & 4 & 83521 & 46 & 3841966 & 70368746289408 & 0.206073 & 0.103037 & 0.20556\\
\bottomrule
\end{tabular}
\caption{Certified explicit dissociated sets $A$ of size $n=ld$ (Theorem~\ref{thm:explicit}). The set $A$ is the lift of
$u=\otimes_k\tilde u_k$; $u$ is determined by $(b,s,Q)$ through Definition~\ref{def:levels}, and the certificate is
$\gamma_1=\dots=\gamma_s=1$. The previous explicit record is $\max A\le0.22002\cdot2^{n}$ \cite{Bo98}.}
\label{tab:main}
\end{table}

Table~\ref{tab:main} is read as follows. In every row the column $f(n)\le$ is within $0.6\%$ of $\Delta_{b,s}$ (the largest
ratio, about $1.006$, occurs at $(11,1)$), so in these examples the finite-scale loss relative to $\Delta_{b,s}$ is small.
The analytic upper bound on the right of Theorem~\ref{thm:explicit} is at least $\Delta_{b,s}>(2/3)^s$. Thus that
upper bound, with $s\le2$, cannot establish $\max A/2^n<\frac12(2/3)^2=2/9$, or improve Bohman's $0.22002$; for $s=3$ one needs
$2^{-b}b^2$ small, and $(7,3)$ still gives $0.2317$ while $(11,3)$ is the first pair below $0.22002$. Increasing $s$
further requires $b$ to grow like $s\log_2s$ to keep $(1+2^{-b})^{(b^s-1)/(b-1)}$ bounded, which is the mechanism behind
Theorem~\ref{thm:rate}; the price is the dimension $d=b^s$ of the vector and the height of the adjugates, about
$3\cdot10^6$ bits for $(17,4)$. The certificate itself is always $s$ adjugates of $b\times b$ integer matrices, and each row
was obtained with at most three trial values of $t$. No row is claimed to be optimal for its size; by
Lemma~\ref{lem:mono}, every row is also an upper bound for $f(n)$ at all larger $n$ (Corollary~\ref{cor:alln}).

\section{Remarks and open problems}\label{sec:remarks}

\subsection{Sizes reached by the construction}
For fixed $(b,s)$, the binary lift produces sizes $n=ld$, where $d=b^s$.
The scale required for a small loss depends on both the perturbation
constant $K_s$ and the spacing of good scalings. More precisely, suppose
that a good scaling can be found in every interval $[X,X+G_{b,s}]$ above
the threshold $8K_s$. Choosing such a scaling with
\[
 2^l+K_s\le Q\le2^l+K_s+G_{b,s}
\]
gives
\[
 \Bigl(\frac Q{2^l}\Bigr)^{d-1}
 \le \exp\!\Bigl(\frac{(d-1)(K_s+G_{b,s})}{2^l}\Bigr).
\]
Thus $d(K_s+G_{b,s})/2^l=o(1)$ suffices to make the finite-scale loss
in Theorem~\ref{thm:explicit} tend to $1$. Positive density of good
scalings in intervals $[X,2X]$ alone gives no such bound on their gaps
near the prescribed point $2^l+K_s$.

Lemma~\ref{lem:AP-heights} proves $K_s=O(b^2d)$ with an absolute
constant. For example, a polynomial bound $G_{b,s}=d^{O(1)}$ permits
$l=O(\log d)$ and hence $n=O(d\log d)$ with a vanishing finite-scale
loss. Such a size estimate requires control of $G_{b,s}$ in addition
to control of $K_s$.

Lemma~\ref{lem:mono} extends any constructed set to every larger size
without changing its normalized maximum. It is useful to compare this
operation with other elementary extensions using their exact ratios.
Write
\[
 \mathcal R(A):=\frac{\max A}{2^{|A|-1}}
\]
for an individual dissociated set $A$.

\paragraph{Adjoining an element larger than all subset sums.}
If $|A|=n$, $S=\sum_{a\in A}a$, and $N=\max A$, then
$A^+=A\cup\{S+1\}$ is dissociated and
\[
 \frac{\mathcal R(A^+)}{\mathcal R(A)}=\frac{S+1}{2N}.
\]
This factor depends on $A$. It is about $n/2$ when $S\sim nN$,
whereas it equals $1$ for $A=\{2^j:0\le j<n\}$.
The identity concerns the two individual sets; it does not assert
that the extremal function $f$ is multiplied by this factor.

\paragraph{Mixed binary lifts.}
Let $u$ have $d$ positive coordinates and be $T$-dissociated, with
$T=2^{l+1}$. Give $r$ coordinates $l+1$ binary digits and the remaining
$d-r$ coordinates $l$ digits, where $1\le r\le d$. The resulting set
$A$ is dissociated, has $ld+r$ elements, and satisfies
\[
 \mathcal R(A)\le
 2^{d-r}\frac{\max u}{T^{d-1}}.
\]
Thus the shape estimate $\max u/T^{d-1}=\Delta(1+o(1))$, when it is
available at this scale, yields the bound
$f(ld+r)\le2^{d-r}\Delta(1+o(1))$. The factor is a bound arising
from this choice of common scale; it need not be the actual ratio of
the mixed lift to a particular binary lift.

\paragraph{Lifts using a general dissociated digit set.}
Let $B\subset\N$ be dissociated, with $m=|B|$, $S=\sum_{a\in B}a$,
and $H=\max B$. Let $u$ be a positive $Q$-dissociated vector of length
$d$, where $Q$ is an integer and $Q\ge S+1$. Then
\[
 A=B\cdot U:=\{a u_i:a\in B,\ 1\le i\le d\}
\]
is dissociated and has $md$ elements: the difference between two digit
sums in each coordinate has absolute value at most $S<Q$, and
dissociation of $B$ handles equality within each coordinate. Its
normalized maximum satisfies the exact identity
\begin{equation}\label{eq:general-digit-lift}
 \mathcal R(A)=
 \frac{H}{2^{m-1}}
 \Bigl(\frac{Q}{2^m}\Bigr)^{d-1}
 \frac{\max u}{Q^{d-1}}.
\end{equation}
The $2^m$ distinct subset sums of $B$ lie among the $S+1$ integers
$0,\dots,S$, so $S+1\ge2^m$. Equality holds precisely for
$B=\{1,2,4,\dots,2^{m-1}\}$: if every integer in $[0,S]$ is a subset
sum, the increasing elements of $B$ must successively be one more
than the sum of their predecessors. For this binary digit set,
choosing $Q=S+1=2^m$ makes the first two factors in
\eqref{eq:general-digit-lift} equal to $1$. For each fixed nonbinary
$B$, the factor $(Q/2^m)^{d-1}$ grows exponentially with $d$ at every
permitted scale $Q\ge S+1$. The additional factor $H/2^{m-1}$ must
also be retained; there is no general loss formula depending only
on $S/2^m$.

\paragraph{Padding the lattice.}
For an admissible pair $(\Lat,v)$ in $\R^d$ of positive height $h$,
set
\[
 \Lat^+=\{(\lambda+kv,-kh):\lambda\in\Lat,\ k\in\Z\}
 \subset V_{d+1},\qquad v^+=e_{d+1}.
\]
This is an admissible pair. Indeed, if the first $d$ coordinates of
$(\lambda+kv,-kh)+tv^+$ lie in $(-1,1)^d$, condition (A2) for the
original pair forces $|k|<1$, hence $k=0$; condition (A1) then gives
$\lambda=0$. The last coordinate consequently gives $|t|<1$.
The projection of $(v,-h)$ onto the orthogonal complement of
$V_d\times\{0\}$ has length $h\sqrt{(d+1)/d}$, so
\[
 \covol(\Lat^+)=\covol(\Lat)h\sqrt{\frac{d+1}{d}},
 \qquad h^+=1,\qquad \Delta^+=h\Delta.
\]
This is an exact factor. It exceeds $1$ for the pairs used here,
whose heights are $(3/2)^s$; an arbitrary admissible pair is only
assumed to have $h>0$.

Monotonicity also shows that the limit of $f(n)$ exists. Every
certified row of Table~\ref{tab:main} bounds $f(n)$ at every larger
size, while an asymptotic rate requires quantitative control of
how the certified sizes grow.

\subsection{The exponent and heights of admissible pairs}
\begin{proposition}\label{prop:height}
Let $(\Lat,v)$ be admissible in $\R^b$, of positive height $h$, and
write $\covol(\Gamma)=1+\eta$. Its $s$-fold tensor power is admissible
in dimension $d=b^s$, with normalized covolume
\[
 \Delta_s=(1+\eta)^{(b^s-1)/(b-1)}h^{-s}.
\]
Conversely, every admissible pair in $\R^b$ satisfies
$\Delta\ge b^{-1/2}$, and therefore
$h\le(1+\eta)\sqrt b$.
\end{proposition}
\begin{proof}
The tensor formula follows by iterating Proposition~\ref{prop:tensor}.
For the second assertion, (A1) and Minkowski's theorem give
\[
 \covol(\Lat)\ge
 2^{-(b-1)}\operatorname{vol}_{b-1}\bigl((-1,1)^b\cap V_b\bigr)
 \ge1,
\]
where the last inequality is Vaaler's cube slicing theorem
\cite{Va79}. Divide by $\sqrt b$ and use
$h\Delta=\covol(\Gamma)=1+\eta$.
\end{proof}

For our base pair, $h=3/2$ and $\eta=2^{-b}$. At the level of lattice
covolumes, the height contributes the factor $h^{-s}$; thus $\log h$
governs this exponential decay in $s$. An improved height can give
an improved bound for $f$, provided the corresponding
primitivisation and lift have suitably controlled size and loss.

More precisely, suppose that a family of base pairs and tensor
powers has $d=b^s\to\infty$,
\[
 h\ge b^\alpha,\qquad
 \eta\frac{b^s-1}{b-1}\le\log2.
\]
Then Proposition~\ref{prop:height} gives
$\Delta_s\le2d^{-\alpha}$. Suppose in addition that effective
primitivisations and binary lifts produce sizes $n=ld$ satisfying
\[
 l=d^{o(1)},\qquad
 f(ld)\le d^{o(1)}\Delta_s.
\]
These two conditions imply
\[
 f(n)\le d^{-\alpha+o(1)}=n^{-\alpha+o(1)}
\]
along the constructed sizes. In the setting of
Theorem~\ref{thm:explicit}, the second condition amounts to
controlling its finite-scale factor
$(Q/2^l)^{d-1}(1+10K_s/Q)$ by $d^{o(1)}$, in addition to satisfying
$2^l\le Q-K_s$.
To obtain the same exponent for all sufficiently large integers,
it is enough that the constructed sizes can be listed increasingly
as $n_j\to\infty$ with
$\log n_{j+1}/\log n_j\to1$; Lemma~\ref{lem:mono} then interpolates
between consecutive sizes. Monotonicity alone does not preserve
this exponent across arbitrarily large multiplicative gaps.
The height bound in Proposition~\ref{prop:height} gives
$\alpha\le1/2$ in this regime, in agreement with the lower bound
$f(n)\gg n^{-1/2}$.

\begin{problem}
Determine
\[
 h^*(b)=\sup\{h:(\Lat,v)\text{ is admissible in }\R^b,
                 \ \covol(\Gamma)\le1+2^{-b}\}.
\]
Does $h^*(b)\to\infty$? Is $h^*(b)\ge b^\alpha$ for some
$\alpha>0$ and all sufficiently large $b$?
\end{problem}

There is a precise packing interpretation with an extension
condition. Put
\[
 C_b=[-1/2,1/2]^b\cap V_b,\qquad
 \nu_b=\operatorname{vol}_{b-1}(C_b).
\]
Condition (A1) says that $\Lat$ is a packing lattice for $C_b$;
its packing density is
\[
 \delta(\Lat)=\frac{\nu_b}{\covol(\Lat)}
             =\frac{\nu_b h}{\sqrt b\,\covol(\Gamma)}.
\]
Restrict to those packing lattices that admit a vector $v$ satisfying
the real-$t$ condition (A2) and
$1\le\covol(\Lat\oplus\Z v)\le1+2^{-b}$. Within this class,
the exact identity
\[
 h=\frac{\sqrt b\,\covol(\Gamma)}{\nu_b}\,\delta(\Lat)
\]
relates large heights to large packing densities, with the
covolume factor between $1$ and $1+2^{-b}$. In particular,
$h\ge b^\alpha$ implies
$\delta(\Lat)\ge \nu_b b^{\alpha-1/2}/(1+2^{-b})$.
A packing lattice satisfying (A1) alone need not supply either the
extension vector or the required full-lattice covolume, so no
equivalence with unrestricted lattice packing is asserted.

\subsection{Siegel's lemma}
As noted in \cite{EP1}, and observed earlier by Schinzel and Aliev \cite{Al08}, for $d\ge2$ let $C_d$ denote the optimal constant in
Siegel's lemma for one equation in $d$ unknowns (so that every nonzero $a\in\Z^d$ has a nonzero integer relation $x$
with $|x|_\infty^{d-1}\le C_d|a|_\infty$; see \cite{BoVa83}). Then $f(d)\ge C_d^{-1}$. Each vector $u$ of Theorem~\ref{thm:explicit} is an explicit integer vector of
height at most $Q^{d-1}\Delta_{b,s}(1+10K_s/Q)$ all of whose nonzero integer relations $c$ satisfy $|c|_\infty\ge Q-K_s$,
so it gives the explicit lower bound $C_d\ge(1-K_s/Q)^{d-1}/\bigl(\Delta_{b,s}(1+10K_s/Q)\bigr)$ for $d=b^s$; for example
$C_{1331}\ge3.15$ and $C_{83521}\ge4.84$.

\subsection*{Acknowledgements}
The underlying lattice construction is due to GPT-6 Astra; we thank Thomas Bloom for its exposition on
\texttt{erdosproblems.com}, and Tom Adamczewski for making the formal proof available. This paper develops the
tensor-preserving perturbation, its structural and local criteria, and the certified constructions described above.
The author discloses that Claude (Anthropic, through the Cursor IDE) and Codex (OpenAI) assisted in developing and
reviewing arguments, implementing and checking computations, and drafting and revising the text. The author takes
full responsibility for the proofs, computations and references. The new results in this paper have not been
formally verified in Lean.
The code and certificate data are available at \url{https://github.com/Chth1z/math-erdos1-dissociated-sets}.

\begin{thebibliography}{99}
\bibitem{Ad26} T.~Adamczewski (repository maintainer). Erd\H{o}s problem \#1: disproof. Lean proof by GPT-6 Astra,
\href{https://github.com/tadamcz/erdos1/tree/0e395153306f34b3829d118b85bdd704136f2843}{\texttt{tadamcz/erdos1}, commit \texttt{0e395153306f}}, 2026.
Primary proof module: \texttt{Erdos1/Resolutions/Erdos1\_219usd\_38h.lean}. Accessed 8 September 2026.
\bibitem{Al08} I.~Aliev. Siegel's lemma and sum-distinct sets. \emph{Discrete Comput.\ Geom.}, 39: 59--66, 2008. \doi{10.1007/s00454-008-9059-9}
\bibitem{EP1} T.~F.~Bloom. Erd\H{o}s Problem \#1, \url{https://www.erdosproblems.com/1}; includes the proof exposition of the construction of GPT-6 Astra (last edited 3 September 2026; status ``disproved, proof verified in Lean''). Accessed 7 September 2026.
\bibitem{Bo96} T.~Bohman. A sum packing problem of Erd\H{o}s and the Conway--Guy sequence. \emph{Proc.\ Amer.\ Math.\ Soc.}, 124: 3627--3636, 1996. \doi{10.1090/S0002-9939-96-03653-2}
\bibitem{Bo98} T.~Bohman. A construction for sets of integers with distinct subset sums. \emph{Electron.\ J.\ Combin.}, 5: Research Paper 3, 1998. \doi{10.37236/1341}
\bibitem{BoVa83} E.~Bombieri and J.~Vaaler. On Siegel's lemma. \emph{Invent.\ Math.}, 73: 11--32, 1983. \doi{10.1007/BF01393823}
\bibitem{CoGu68} J.~H.~Conway and R.~K.~Guy. Sets of natural numbers with distinct sums. \emph{Notices Amer.\ Math.\ Soc.}, 15: 345, 1968. See also R.~K.~Guy, \emph{Unsolved Problems in Number Theory}, 3rd ed., Springer, 2004, Problem C8.
\bibitem{DFX21} Q.~Dubroff, J.~Fox, and M.~W.~Xu. A note on the Erd\H{o}s distinct subset sums problem. \emph{SIAM J.\ Discrete Math.}, 35: 322--324, 2021. \doi{10.1137/20M1385883}
\bibitem{El86} N.~D.~Elkies. An improved lower bound on the greatest element of a sum-distinct set of fixed order. \emph{J.\ Combin.\ Theory Ser.\ A}, 41: 89--94, 1986.\newline\doi{10.1016/0097-3165(86)90116-0}
\bibitem{Er56} P.~Erd\H{o}s. Problems and results in additive number theory. In \emph{Colloque sur la Th\'eorie des Nombres, Bruxelles, 1955}, pages 127--137. Georges Thone, Li\`ege; Masson, Paris, 1956.
\bibitem{Gu82} R.~K.~Guy. Sets of integers whose subsets have distinct sums. In \emph{Theory and Practice of Combinatorics}, North-Holland Math.\ Stud.\ 60 / Ann.\ Discrete Math.\ 12, pages 141--154. North-Holland, Amsterdam, 1982. \doi{10.1016/S0304-0208(08)73500-X}
\bibitem{HoKa23} T.~Horesh and Y.~Karasik. Equidistribution of primitive lattices in $\R^n$. \emph{Q.\ J.\ Math.}, 74: 1253--1294, 2023. \doi{10.1093/qmath/haad008}
\bibitem{Lu88} W.~F.~Lunnon. Integer sets with distinct subset-sums. \emph{Math.\ Comp.}, 50: 297--320, 1988. \doi{10.1090/S0025-5718-1988-0917837-5}
\bibitem{St23} S.~Steinerberger. Some remarks on the Erd\H{o}s distinct subset sums problem. \emph{Int.\ J.\ Number Theory}, 19: 1783--1800, 2023. \doi{10.1142/S1793042123500860}
\bibitem{Va79} J.~D.~Vaaler. A geometric inequality with applications to linear forms. \emph{Pacific J.\ Math.}, 83: 543--553, 1979. \doi{10.2140/pjm.1979.83.543}
\end{thebibliography}

\end{document}
